We shall be using the following generalised notion of orbits for a subset of a nominal set.

\begin{definition}
  Let $\Xcal$ be a nominal set and $S\subset\Xcal$ have finite support. For each finite $N\subseteq\Atoms$ we set:\
$
\orb[N](S) = \{ \orbs[N](x)\mid x\in S\} 
$.
\end{definition}

\subsection{Proof of \cref{lem:negfreesize}}
%\begin{lemma}\label{lem:negfreesize}
\negfreesize*
%	For any formula $\phi$, $|!(\phi)|$ is bounded by $|\phi| - n$, where $n$ is the number of negations in $\phi$. 
%\end{lemma}

\begin{proof}
	We perform an induction on the measure $(m, n)$ of $\phi$, where $m$ is the number of non-negation connectives and constructs that appear in $\phi$, 
	with the following base cases:
 	\begin{itemize}
		\item $|!(u = v)| = |u = v| = 2$,
		\item $|!(u \neq v)| = |u \neq v| = 2$,
		\item $X(\vec u) = {!(X(\vec u))}$, hence they have the same size.
	\end{itemize}
	Next, we perform inductive steps.
	The cases where $\phi$ is not of the form $\lnot \psi$ follow from the inductive hypothesis. 
	For instance, if $\phi = \phi_1 \lor \phi_2$, then 
	\[
		|\phi| = |\phi_1| + |\phi_2| + 1 \overset{\text{IH}}{\le} |!(\phi_1)| + |!(\phi_2)| + 1 = |!(\phi)|
	\]
	Thus, we only examine the more interesting cases where $\phi$ is of the form $\lnot \psi$. 
	We perform a case analysis on $\phi$:
	    \begin{itemize}
		\item $\lnot (u = v)$: this formula has $n = 1$, and $|\lnot (u = v)| = |\{u,v\}|$. 
		Then, we have
		$|!(\lnot(u = v))| = |u \neq v| = |\lnot (u = v)| - n = |\{u,v\}|$.
		Similar can be shown for $\lnot (u \neq v)$.
		%
		\item $\lnot (\phi_1 \lor \phi_2)$: this formula has size $2 + |\phi_1| + |\phi_2|$. 
		By definition, we have that $|!(\lnot (\phi_1 \lor \phi_2))| = |!(\lnot \phi_1) \land !(\lnot \phi_2)| = |!(\lnot \phi_1)| + |!(\lnot \phi_2)| + 1$. 
		Let us write $n_1$ for the number of negations in $\phi_1$, and $n_2$ for the number of negations in $\phi_2$. 
		Thus, the number of negations in $\phi$ would be $n = n_1 + n_2 + 1$.
		The inductive hypothesis can be used as each $\lnot \phi_i$ will have one less connective than $\phi$, that being the $\lor$, meaning
		we have that $|!(\lnot \phi_i)| \le |\lnot \phi_i| - n_i - 1$ for $i \in \{1,2\}$.
		Hence, we have:
		\begin{align*}
			|!(\lnot \phi_1)| +| !(\lnot \phi_2)| + 1 & \le (|\lnot\phi_1| -n_1-1) + (|\lnot\phi_2| -n_2- 1) + 1\\
			&= |\phi_1| + |\phi_2| -(n_1+n_2+1)+ 2\\
			&= (|\phi_1| + |\phi_2| + 2) - n
		\end{align*}
		Similar can be shown for $\lnot (\phi_1 \land \phi_2)$.
		%
		\item $\lnot \bigvee\nolimits_x \phi'$: this formula has size $2 + |\phi'|$. 
		By definition, we have that $!(\lnot \bigvee\nolimits_x \phi') = \bigwedge\nolimits_x !(\lnot \phi')$. 
		The inductive hypothesis can be used as $\lnot \phi'$ has one less connective than $\phi$, that being $\bigvee\nolimits_x$, meaning
		we have that $|!(\lnot \phi')| \le |\lnot \phi'| - n$.
		Hence, we have:
		\begin{align*}
			|!(\lnot \phi')| + 1 &\le (|\lnot \phi'| - n) + 1= |\phi'| - n + 2
		\end{align*}    
		A similar argument can be shown for $\lnot\bigwedge\nolimits_x \phi'$.
		%
        \item $\lnot \fresh{x}\phi'$ this formula has size $2 + |\phi'|$. 
		By definition, we have that $!(\lnot \fresh{x} \phi') = \fresh{x}!(\lnot \phi')$. 
		The inductive hypothesis can be used as $\lnot \phi'$ has one less connective than $\phi$, that being $\fresh{x}$, meaning
		we have that $|!(\lnot \phi')| \le |\lnot \phi'| - n$.
		Hence, we have:
		\begin{align*}
			|!(\lnot \phi')| + 1 &\le (|\lnot \phi'| - n) + 1= |\phi'| - n + 2
		\end{align*}
		%
		\item $\lnot \diam{t, \vec u} \phi'$: 
		this formula has size $2 + |\vec u| + |\phi'|$. 
		By definition, we have that $!(\lnot\diam{t, \vec u}\phi') = \sq{t, \vec u}!(\lnot \phi')$.
		The inductive hypothesis can be used as $\lnot \phi'$ has one less connective than $\phi$, that being $\diam{t, \vec u}$, meaning we have that $|!(\lnot \phi')| \le |\lnot \phi'| - n$. 
		Hence, we have:
		\[
		|!(\lnot \phi')| + 1 + |\vec u| \le (|\lnot \phi'| - n) + 1 + |\vec u| = 2 + |\phi'| + |\vec u| - n 
		\]
		A similar argument can be shown for $\lnot\sq{t, \vec u}\phi'$.
		% 
		\item $\lnot (\mu X(\vec x).\phi')(\vec u)$: this formula has size $2 + |\phi'| + |\vec u|$. 
		By definition, we have that $!(\lnot (\mu X(\vec x).\phi')(\vec u)) = (\nu X(\vec x).!(\lnot \phi'_X))(\vec u)$.
		The inductive hypothesis can be used as $\lnot \phi'_X$ has one less construct than $\phi$, that being $\mu X$, meaning
		we have that $|!(\lnot\phi'_X)| \leq |\lnot \phi'_X| - n'$ where $n'$ is the number of negations in $\lnot \phi'_X$.
		This means that $|\lnot \phi'_X| - n' = |\phi'| - n$ as both have the number of negations removed,
		hence, we have:
		\begin{align*}
			1 + |!(\lnot \phi'_X)| + |\vec u| &\le 1 + |\lnot \phi'_X| - n' + |\vec u| \\
			&= 1 + |\phi'| - n + |\vec u|\\
			&\le 2 + |\phi'| + |\vec u| - n
		\end{align*}
		A similar argument can be shown for $\lnot (\nu X(\vec x).\phi')(\vec u)$.
		%
		\item $\lnot \lnot \phi'$: this formula has size $2 + |\phi'|$. 
		By definition, we have that $!(\lnot \lnot \phi') = !(\phi')$. 
		The inductive hypothesis can be used as the number of connectives for $\lnot \lnot \phi'$ and $\phi'$ are the same, however, there are two less negations. 
		Hence, by the inductive hypothesis we have $|!(\phi')| \le |\phi'| - (n - 2)$, 
		meaning $|!(\phi)| = |!(\phi')| \le |\phi'| + 2 - n $.
%		:
%		\begin{align*}
%			|!(\phi)| = |!(\phi')| &\le |\phi'| + 2 - n 
%		\end{align*}
	\end{itemize}
%    \begin{itemize}
%    
%    \end{itemize}
\end{proof}

\subsection{Proof of  \cref{lem:freshClosProps}}

\begin{lemma}\label{lem:unionOrb2}
	Given two sets $X_1,X_2\subseteq \Xcal$ such that $\supp(X_i) \subseteq \S \subseteq \Atoms$ for $i \in \{1, 2\}$,
	then $\orb[\S](X_1 \cup X_2) = \orb[\S](X_1) \cup \orb[\S](X_2)$.
\end{lemma}

\begin{proof}
	By definition, we have that:
	\begin{align*}
		\orb[\S](X_1 \cup X_2) &= \{\orbs[\S](x) \mid x \in X_1 \cup X_2\} 
		\\ &= \{\orbs[\S](x) \mid x \in X_1\} \cup \{\orbs[\S](x) \mid x \in X_2\} 
		\\ &= \orb[\S](X_1) \cup \orb[\S](X_2)    
	\end{align*}
\end{proof}

\begin{lemma}\label{lem:unionOrb3} 
	Let $\S \subseteq \Atoms$ such that $\S$ is finite and $\Xcal$ be a nominal set.
	Then for any $X \subseteq \Xcal$ and $b \notin \S$ such that $\supp(X) \subseteq \S \cup \{b\}$, 
	there exists a 1-1 function 
	$f: \orb[\S](\bigcup_{a \notin \S} (a\ b)\cdot X) \rightarrow \orb[\S \cup \{b\}](X)$.
\end{lemma}
\begin{proof}
	We define the function $f$ as follows: for each $O \in \orb[\S](\bigcup_{a \notin \S} (a\ b)\cdot X)$, such that $O = \orbs[\S]((a\ b) \cdot x)$ for some $x \in X$ and some $a \notin \S$, then $f(O) = \orbs[\S \cup \{b\}](x)$.
	The fact that $f$ is a function follows from the fact that the orbits in the codomain are taken over a larger set than the orbits in the domain, it remains to show $f$ is $1-1$. 
	Let us choose some $x_1, x_2 \in X$, so $f(\orbs[\S]((a\ b) \cdot x_i)) = \orbs[\S \cup \{b\}](x_i)$. 
	Assume $\orbs[\S]((a\ b) \cdot x_1) \neq \orbs[\S]((a\ b)\cdot x_2)$, meaning there is no $\pi \in \fix(\S)$ such that $\pi \cdot ((a\ b)\cdot x_1) = (a\ b)\cdot x_2$.     
	Let us also assume $f(\orbs[\S]((a\ b)\cdot x_1)) = f(\orbs[\S]((a\ b)\cdot x_2)) = \orbs[\S \cup \{b\}](x_1)$. 
	This means there must be some $\pi' \in \fix(\S \cup \{b\})$ such that $\pi' \cdot x_1 = x_2$, and so we also have $((a\ b)\circ \pi') \cdot x_1 = (a\ b) \cdot x_2$.
	Let us write $\hat \pi = (a\ b)\circ \pi' \circ (a\ b)$, so $((a\ b)\circ \pi') \cdot x_1 = \hat \pi \cdot ((a\ b)\cdot x_1) = (a\ b)\cdot x_2$. We have that $\hat \pi \in \fix(\S)$, however, this is a contradiction as this would mean $\orbs[\S]((a\ b) \cdot x_1) = \orbs[\S]((a\ b)\cdot x_2)$. 
	Hence, $f$ is $1-1$.
\end{proof}

\begin{lemma}\label{lem:freshClosProps1}
		Let $\Lcal$ be a nominal LTS.
	Given a negation-free formula $\phi_0$ with set of free recursion variables $\cal Z$ and a configuration $(s_0,H_0)$, thenfor all $(\phi,\gamma,\theta)\in\clos(\phi_0)$, the free variables of $\phi$ are included in $\dom(\theta)\uplus \cal Z$ and all formulas in $\rng(\theta)$ are fixpoints and have free variables in $\dom(\theta)\uplus\Zcal$.
\end{lemma}

\begin{proof}
	We denote $\fv(\psi)$ for the free variables of a formula $\psi$, 
	and we do induction on the derivation of $(\phi,\gamma,\theta)$, starting from $(\phi_0,\{\}, \{\})$. 
	The base cases are clear; let $(\phi',\gamma',\theta')$ be the parent of $(\phi, \gamma, \theta)$. 
	We do a case analysis on $\phi'$ as follows:
	%\ntnote{unclear what the case analysis below is on}
	\begin{itemize}
		\item $\phi_1 \lor \phi_2$: By the inductive hypothesis,  
		we have that $\fv(\phi_1), \fv(\phi_2) \subseteq \fv(\phi_1 \lor \phi_2) \subseteq \dom(\theta)\uplus \Zcal$.
		A similar argument can be shown for $(\phi_1 \land \phi_2)$.
		%
		\item $\diam{\ell}\phi$: By the inductive hypothesis, 
		we have that $\fv(\phi') = \fv(\phi) \subseteq \dom(\theta) \uplus \Zcal$. 
		A similar argument can be shown for $\sq{\ell}\phi, \bigvee_x \phi'',  \bigwedge_x \phi''$ and $\fresh{x}\phi''$.
		%
		\item $\fresh{x}\phi''$: there exists an $a \in \Atoms$ such that $\phi = \phi''\{a/x\}$. By the inductive hypothesis, we have that $\fv(\phi) = \fv(\phi'') = \fv(\phi') \subseteq \dom(\theta) \uplus \Zcal$.			
		%
		\item $(\mu X(\vec x).\phi'')(\vec a)$: $X$ will become free in $\phi=\phi''\{\vec c/\vec x\}$, for some $\vec c$, and $\theta=\theta'[(\mu X(\vec x).\phi'')/X]$.
		By definition, $X\notin\Zcal$.
		By the inductive hypothesis, we have that 
		$\fv(\phi) =\fv(\phi'') \cup \{X\} \subseteq (\dom(\theta') \uplus \Zcal)\cup \{X\}=\dom(\theta) \uplus \Zcal$. 
		By the inductive hypothesis, the new elements in the range of $\theta$ have free variables in $\Zcal$
		and are fixpoints.
		A similar argument can be shown for $(\nu X(\vec x).\phi'')(\vec a)$.
	\end{itemize}
\end{proof}

\freshClosProps*


\begin{proof}
	\begin{enumerate}
		%
		\item 
		For $\phi$ to be in $\clos^*(\phi_0)$, it suffices to show that there exists some $(\phi', \gamma', \theta') \in \clos(\phi_0)$ such that $\phi = \phi'\{\theta'\}\{\gamma'\}$. 
		We perform induction on the path leading to $(\phi', \gamma', \theta')$ starting from $(\phi_0, \{\}, \{\})$ with the following base case:
		\begin{itemize}
			\item $(s, H, \xi, \phi) \in \Gcal(\Lcal, \phi_0,\xi_0, s_0)$ implies $\phi \in \clos(\phi_0)$ and so $(\phi', \gamma', \theta') \in \clos(\phi_0)$.
		\end{itemize}
		Next, we perform inductive steps. Suppose the parent position of $(s, H,\xi,\phi)$ has formula:
		\begin{itemize}
			\item $\phi_1 \lor \phi_2$: by the inductive hypothesis, there is a triple $(\psi,\gamma,\theta)\in\clos(\phi_0)$ such that $\phi_1\lor\phi_2=\psi\{\theta\}\{\gamma\}$. From Lemma~\ref{lem:freshClosProps1}, there is $\psi=\psi_1\lor\psi_2$ and therefore $\phi=\psi_1\{\theta\}\{\gamma\}\lor\psi_2\{\theta\}\{\gamma\}$, and $(\psi_1,\gamma,\theta),(\psi_2,\gamma,\theta)\in\clos(\phi)$, as required. The case for  $\phi_1 \land \phi_2$ is similar.
			%			
			\item $\diam{\ell}\phi''$: 
			by the inductive hypothesis, there is a triple $(\psi,\gamma, \theta)\in\clos(\phi_0)$ such that $\diam{\ell}\phi''=\psi\{\theta\}\{\gamma\}$.
			From Lemma~\ref{lem:freshClosProps1}, there is $\psi=\diam{\ell}\psi''$, and $(\psi'', \gamma,\theta)\in \clos(\phi_0)$, as required. 
			The case for $\sq{\ell}\phi''$ is similar.
			%			
			\item $\fresh{x}\phi''$:
			by the inductive hypothesis, there is a triple $(\psi, \gamma, \theta) \in \clos(\phi_0)$ such that $\fresh{x}\phi''=\psi\{\theta\}\{\gamma\}$.
			From Lemma~\ref{lem:freshClosProps1}, there is $\psi=\fresh{x}\psi''$, and $(\psi'', (\gamma, \{a/x\}), \theta) \in \clos(\phi_0)$ as required.
			The cases for $\bigvee_x \phi''$ and $\bigwedge_x \phi''$ are similar. 
			%			
			\item $(\mu X(\vec x).\phi'')(\vec a)$: by the inductive hypothesis, there is a triple $(\psi, \gamma,\theta) \in \clos(\phi_0)$ such that $\psi\{\theta\}\{\gamma\} = (\mu X(\vec x).\phi'')(\vec a)$. 
			If $\psi$ is a variable, then $\psi = X(\vec a)$ and $(\mu X(\vec x).\phi'')(\vec a) = \theta(X(\vec a))\{\gamma\}$. 
			By Definition~\ref{def:closFresh}, there must be (for some $\vec b$) $((\mu X(\vec x).\phi')(\vec b),\gamma, \theta) \in \clos(\phi_0)$, meaning $(\phi'',(\gamma, \{\vec a / \vec x\}), \theta) \in \clos(\phi_0)$ and $\phi''\{\theta\}\{\gamma\}\{\vec a / \vec x\} = (\phi''\{\mu X(\vec x) / X\})\{\gamma\}\{\vec a / \vec x\}$ as required.
%			 \hbnote{check this again}
			\\
			If $\psi$ is not a variable, then $\psi = (\mu X(\vec x).\hat \phi)(\vec a))$ such that $ (\mu X(\vec x).\hat \phi)(\vec a))\{\theta\}\{\gamma\} =  (\mu X(\vec x).\phi'')(\vec a))$. 
			Let $\theta'' = \theta[\mu X(\vec x).\phi' / X]$.Then, from $(\psi,\gamma, \theta) \in \clos(\phi_0)$ we have $(\phi''\{\vec a/ \vec x\}, \gamma,\theta') \in \clos(\phi_0)$ and $(\phi'\{\vec a/ \vec x\})\{\theta'\}\{\gamma\} = (\phi'\{\theta'\})\{\gamma\}\{\vec a/ \vec x\} = (\phi'\{\mu X(\vec x).\phi' / X\})\{\gamma\}\{\vec a / \vec x\}$ as required.
%			\hbnote{check this again}
			
			The case for $(\nu X(\vec x).\phi')(\vec a)$ is similar.
		\end{itemize}		
		\item 
		% Let us write $\orb[\phi',\theta']$ to be the orbit that fixes the support of all formulas that appear in the path from $\clos(\phi, [])$ to $\clos(\phi', \theta')$.
		%\ntnote{update this with $\psi,\theta$. Actually, this does not work bc later we need a version of Lemma~11 that does not hold. Instead, try taking orbits wrt set $N$ which is the support of all formulas from the root $(\phi,[])$ to $(\psi,\theta)$}
		%\fixednote{better write $\orb[\phi'] = \orb[\supp(\phi,\phi')]$ as $\phi$ may contain names that do not appear in $\phi'$ anymore, but may appear in $\theta$}.
%		By definition, we have that $|\orb[\supp(\phi)](\clos(\phi))| = |\orb[\supp(\phi)](\clos(\phi))|$,
		%\fixednote{orbits needed}, 
%		so it remains to show that $|\orb[\supp(\phi)](\clos(\phi))| \le |\phi|\cdot m(\phi)$.
		%\fixednote{update}
		Let $(\psi, \gamma, \theta) \in \clos(\phi)$. 
		We write $\clos_\phi(\psi, \gamma, \theta)$ for the subset of $\clos(\phi)$ derived from and including $(\psi,\gamma, \theta)$, and let $\Ncal$ be the support of all formulas and name substitutions in the path from $(\phi,\{\}, \{\})$ to $(\psi,\gamma, \theta)$,
		therefore $\supp(\phi, \psi, \gamma, \theta) \subseteq \Ncal$.
		%\fixednote{maybe add that therefore $\supp(\phi,\psi,\theta)\subseteq\Ncal$ (this answers a comment later on as well)}
		Without loss of generality, for the fixpoint case we consider a more general $\clos$ rule for this case, which goes from $((\mu X(\vec x).\psi')(\vec a),\gamma, \theta)$ to $(\bigvee\nolimits_{x_1}\dots\bigvee\nolimits_{x_n}\psi',\gamma',\theta')$ where $\theta' = \theta\{\mu X(\vec x).\psi' / X\}$ and $\gamma' = (\gamma, \{\vec a / \vec x\})$. 
		This change allows for each $x \in \vec x$ to become free sequentially, in contrast to making all $x \in \vec x$ free at once. 
		The impact on result is unaffected as this will add more intermediate formulas (i.e., each $\bigvee_{x_i}$).
		% \fixednote{expland a bit more on "goes to", also, you need $\theta'$}
		Next we define the following function:
		\[
		m'(\phi', \gamma', \theta') = \frac{(\|\phi'\| + |\Ncal|)!}{|\Ncal|!}
		\]
		noting that $m(\phi) = m'(\phi, [])$ (where it does not cause ambiguity, the $\theta'$ and $\gamma'$ parameters in $m'$ may be dropped for economy).
		We show that $|\orb[\Ncal](\clos_\phi(\psi, \gamma, \theta))| \le |\psi|\cdot m'(\psi, \gamma, \theta)$
		%\fixednote{orbit subscript} 
		by performing induction on $\psi$:
		%\ntnote{the m needs $\theta$ as a second argument; maybe you can say we drop it for economy when not used}
		\begin{itemize}
			\item $|\orb[\Ncal](\clos_\phi(a = a, \gamma, \theta))|= 1 \le 2 = |\psi|\cdot m'(\psi)$.
			Similar can be shown for $a \neq a$.
			\item $|\orb[\Ncal](\clos_\phi(a = b, \gamma, \theta))| = 1 \le 2 = |\psi|\cdot m'(\psi)$.
			Similar can be shown for $a \neq b$.
			\item $|\orb[\Ncal](\clos(X(\vec a), \gamma, \theta))| = 1 \le 1 + |\vec a| = |\psi|\cdot m'(\psi)$.
		\end{itemize}
		Next we perform inductive steps:
		\begin{itemize}
			\item $\psi_1 \lor \psi_2$: by the inductive hypothesis, $|\orb[\Ncal](\clos_\phi(\psi_i,\gamma, \theta))| \le |\psi_i|\cdot m'(\psi_i)$ for $i \in \{1, 2\}$.
			By Lemma~\ref{lem:unionOrb2},
			%\fixednote{I think the lemma (\ref{lem:unionOrb}) cannot be applied in this direction?} 
			we have that $|\orb[\Ncal](\clos_\phi(\psi_1 \lor \psi_2,\gamma, \theta))| \le |\orb[\Ncal](\clos_\phi(\psi_1, \gamma, \theta))| + |\orb[\Ncal](\clos_\phi(\psi_2,\gamma, \theta))|$.
			As, for $i \in \{1,2\}$, $\|\psi_i\| \le \|\psi\|$ we have $m'(\psi_i) \le m'(\psi)$. 
			Thus we have:
			\begin{align*}
				|\orb[\Ncal](\clos_\phi(\psi_1 \lor \psi_2,\gamma, \theta))| &\le |\orb[\Ncal](\clos_\phi(\psi_1, \gamma, \theta))| + |\orb[\Ncal](\clos_\phi(\psi_2,\gamma, \theta))| 
				\\&\le (|\psi_1|\cdot m'(\psi_1)) + (|\psi_2|\cdot m'(\psi_2)) \le |\psi|\cdot m(\psi)
			\end{align*}%\fixednote{a m missing at end}
			A similar argument can be shown for $\psi_1 \land \psi_2$.
			% 
			\item $\diam{\ell}\psi'$: by the inductive hypothesis, $|\orb[\Ncal](\clos_\phi(\psi', \gamma, \theta))| \le |\psi'|\cdot m'(\psi')$.
			By definition
			we have that $|\orb[\Ncal](\clos_\phi(\psi, \gamma, \theta))| \le |\orb[\Ncal](\clos_\phi(\psi',\gamma,\theta))| + 1$ and $m'(\psi') = m'(\psi)$.
			%\fixednote{the ineq is the other way round than used below. Maybe find an expression with equality instead}
			Thus we have:
			\begin{align*}
				|\orb[\Ncal](\clos_\phi(\psi,\gamma,\theta))| &\le |\orb[\Ncal](\clos_\phi(\psi',\gamma,\theta))| + 1
				\\&\le |\psi'|\cdot m'(\psi') + 1 
				\\&\le |\psi|\cdot m'(\psi)
			\end{align*}
			A similar argument can be shown for $\sq{\ell}\psi'$.
			%
			\item $\bigvee_x \psi'$: by the inductive hypothesis, $|\orb[\Ncal \cup \{a\}](\clos_\phi(\psi',(\gamma, \{a/x\}), \theta))| \le |\psi'|\cdot m'(\psi',(\gamma, \{a/x\}),\theta)$ for all $a \in \Atoms$.
			By definition, we have that 
			%\ntnote{here we should be counting names in $\Ncal$ and not in $\Ncal$ (bc of the orbits mentioning $\Ncal$); therefore, for the induction we should probably use a larger m, namely $m'(\phi',\theta)=(\|\phi'\|+|\Ncal|)!/|\Ncal|!$}
			\begin{align*}
				\clos_\phi(\psi,\gamma, \theta) &= (\bigcup\nolimits_{a \in \Ncal} \clos_\phi(\psi',(\gamma, \{a/x\}), \theta))\\ &\cup (\bigcup\nolimits_{a \notin \Ncal} \clos_\phi(\psi',(\gamma, \{a/x\}), \theta)) \cup \{(\psi,\gamma, \theta)\}
			\end{align*}
			we know, for any $b \notin \Ncal$:
			%\fixednote{here you need $b\notin N$}
			\[
			\orb[\Ncal](\bigcup\nolimits_{a \notin \Ncal} \clos_\phi(\psi',(\gamma, \{a/x\}), \theta) = \orb[\Ncal](\bigcup\nolimits_{a \notin \Ncal} (a\ b) \cdot \clos_\phi(\psi',(\gamma, \{\vec a / \vec x\}), \theta))
			\]
			%
			and then by Lemma~\ref{lem:unionOrb3} we obtain 
			\[
			|\orb[\Ncal](\bigcup\nolimits_{a \notin \Ncal} (a\ b) \cdot \clos_\phi(\psi',(\gamma, \{\vec a / \vec x\}), \theta))| \le |\orb[\Ncal \cup \{b\}](\clos_\phi(\psi',(\gamma, \{\vec a / \vec x\}), \theta))|
			\]
			As each such, $\clos_\phi(\psi',(\gamma, \{a/x\}), \theta)$ and $\bigcup\nolimits_{a \notin \Ncal} (a\ b) \cdot\clos_\phi(\psi',(\gamma, \{ a /  x\}), \theta)$ are closed under permutations fixing $\supp(\psi,\theta)$, by Lemma~\ref{lem:unionOrb2} for $b \notin \Ncal$ we have that 
			%\fixednote{$b\notin N$}
			\begin{align*}
				\orb[\Ncal](\clos_\phi(\psi, \theta)) 
				&= \bigcup\nolimits_{a \in \Ncal} \orb[\Ncal](\clos_\phi(\psi',(\gamma, \{a/x\}), \theta)) 
				\\
				%\cup \orb[\psi](\clos_\phi(\psi'\{b/x\}, \theta))\\
				&\quad\cup \orb[\Ncal]\left(\bigcup\nolimits_{a \notin \Ncal} (a\ b) \cdot\clos_\phi(\psi',(\gamma, \{a/x\}), \theta)\right)
				\cup\{(\psi, \gamma,\theta)\} 
			\end{align*}
			and hence
			\begin{align*}
				|\orb[\Ncal](\clos_\phi(\psi, \gamma,\theta)) |
				% &= \orb[\psi](\bigcup\nolimits_{a \in \supp(\psi)} \clos_\phi(\psi'\{a/x\}, \theta) \cup \clos_\phi(\psi'\{b/x\}, \theta) \cup (\psi, \theta))\\
				%&= \orb[\psi](\bigcup\nolimits_{a \in \supp(\psi,\theta)} \clos_\phi(\psi'\{a/x\}, \theta)) \cup \orb[\psi](\clos_\phi(\psi'\{b/x\}, \theta)) \cup \orb[\psi](\{\psi, \theta\}) 
				%\\
				&\leq \sum\nolimits_{a \in \Ncal} |\orb[\Ncal](\clos_\phi(\psi',(\gamma, \{a/x\}), \theta))| \\
				&\quad + |\orb[\Ncal \cup \{b\}](\clos_\phi(\psi',(\gamma, \{ a /  x\}), \theta))|+1\\
				&\overset{\text{IH}}{\leq}
				\sum\nolimits_{a \in \Ncal} |\psi'|\cdot m'(\psi',(\gamma, \{a/x\}),\theta) \\
				&\quad + |\psi'|\cdot m'(\psi',(\gamma, \{ a /  x\})\theta)+1
			\end{align*}
			For $a \in \Ncal$, then $m'(\psi',(\gamma, \{a/x\}),\theta) = (|\Ncal| + \|\psi\| - 1)!/|\Ncal|! = m'(\psi,\gamma,\theta)/(|\Ncal|+\|\psi\|)$. 
			For $b \notin \supp(\psi',\gamma,\theta)$, then $m'(\psi',(\gamma, \{\vec a / \vec x\}),\theta) = (|\Ncal| + \|\psi\|) / (|\Ncal| + 1)! = m'(\psi)/(|\Ncal| + 1)$.
			Thus, we obtain:
			\begin{align*}
				|\orb[\Ncal](\clos_\phi(\psi,\gamma, \theta)| 
				&\le 
				\sum\nolimits_{a \in \Ncal} |\psi'|\cdot m'(\psi,(\gamma, \{a/x\}),\theta)/(|\Ncal|+1)\\
				&\quad + |\psi'|\cdot m'(\psi,\gamma,\theta)/(|\Ncal|+1)+1
				\\
				&=|\psi'|\cdot m'(\psi,\gamma,\theta)/(|\Ncal|+1)\cdot (|\Ncal|+1)+1 
				\\
				&=|\psi'|\cdot m'(\psi,\gamma,\theta)+1
				\leq |\psi|\cdot m'(\psi,\gamma,\theta)
			\end{align*}
			A similar argument can be shown for $\bigwedge\nolimits_x \psi'$ and $\fresh{x} \psi'$.
			%
			\item $(\mu X(\vec x).\psi')(\vec a)$: 
			Let $\theta' = \theta[\mu X(\vec x).\psi' / X]$.
			%\fixednote{technically speaking, we cannot exactly apply the IH here as $\hat\psi$ is not in the closure and therefore it does not have a $\Ncal$. W can say at the beginning of the proof for part 3 that we consider a more general clos rule for this case, which takes us from $(\phi,\theta)$ to $(\hat\psi,\theta')$}
			%
			Let us write $\hat \psi = \bigvee\nolimits_{x_1}\dots\bigvee\nolimits_{x_n}\psi'$ where $x_i$ is the $i$th element of $\vec x$.
			By definition, we have that
			\[
			\clos_\phi((\mu X(\vec x).\psi')(\vec a), \gamma, \theta) = \{(\mu X(\vec x).\psi')(\vec a), \gamma, \theta)\} \cup \clos_\phi(\hat \psi, (\gamma, \{\vec a / \vec x\}), \theta')
			\]
			and by Lemma~\ref{lem:unionOrb2}:
			% noting $\supp(\psi,\theta) \supseteq\supp(\hat \psi,\theta')$:
			\[
			\orb[\Ncal](\clos_\phi((\mu X(\vec x).\psi')(\vec a),\gamma, \theta)) \subseteq \{((\mu X(\vec x).\psi')(\vec a),\gamma, \theta)\} \cup \orb[\Ncal](\clos_\phi(\hat \psi,(\gamma, \{\vec a/\vec x\}), \theta'))
			\]
			By the inductive hypothesis, we have:
			\[
			|\orb[\Ncal](\clos_\phi(\hat \psi,(\gamma, \{\vec a/\vec x\}) \theta'))| \le |\hat \psi|\cdot m'   (\hat \psi,(\gamma, \{\vec a/\vec x\}),\theta')
			\]
			We note that
			$m'(\hat \psi,(\gamma, \{\vec a/\vec x\}),\theta') = m'(\psi,\gamma,\theta)$, thus obtaining:
			%\ntnote{I think Lemma 11 is applied the wrong way round here?}
			\begin{align*}
				|\orb[\Ncal](\clos_\phi(\psi,\gamma, \theta))| &\le |\orb[\Ncal](\clos_\phi(\hat \psi,(\gamma, \{\vec a/\vec x\})\theta'))| + 1
				\\ &\overset{\text{IH}}{\le} |\hat \psi|\cdot m'(\hat \psi) + 1
				\\ &\leq (|\vec x| + |\psi'|)\cdot m'(\psi,\gamma,\theta) + 1
				\leq |\psi|\cdot m'(\psi,\gamma,\theta)
			\end{align*}
			A similar argument can be shown for $(\nu X(\vec x).\psi')(\vec a)$.
		\end{itemize}
		\item    We show that for any $(\phi',\gamma, \theta) \in \clos(\phi)$, $|\supp(\phi')| + \|\phi'\| \le |\supp(\phi)| + \|\phi\|$.
		% We show that for any $(\phi', \theta) \in \clos(\phi)$ with parent $(\phi'', \theta)$, then $|\supp(\phi')| \le |\supp(\phi'')| + \|\phi''\|$.
		Let $(\phi'',\gamma'\theta')$ be the parent node of $(\phi',\gamma' \theta)$.
		We do induction on the derivation of $(\phi',\gamma, \theta)$, starting from $(\phi,\{\}, \{\})$; the base cases are clear.
		We perform a case analysis on $\phi''$ as follows:
		\begin{itemize}
			\item $\phi_1 \lor \phi_2$: 
			By the inductive hypothesis, $|\supp(\phi_1 \lor \phi_2)| + \|\phi_1 \lor \phi_2\| \le |\supp(\phi)| + \|\phi\|$. 
			We have that $\supp(\phi_1,\phi_2) \subseteq \supp(\phi_1 \lor \phi_2)$ and $\|\phi_i\| \le \|\phi_1 \lor \phi_2\|$ for $i \in \{1,2\}$, and so we obtain:
			\[
			|\supp(\phi_i)| \le |\supp(\phi_1 \lor \phi_2)| \le  |\supp(\phi)| + \|\phi\|
			\]
			Similar can be shown for $\phi_1 \land \phi_2$.
			%
			\item $\diam{\ell} \phi'$: By the inductive hypothesis, $|\supp(\diam{\ell}\phi')| + \|\phi'\| \le |\supp(\phi)| + \|\phi\|$, and by definition, $\supp(\phi')\subseteq \supp(\diam{\ell}\phi')$ and $\|\phi'\| = \|\diam{\ell}\phi'\|$, therefore:
			\[
			|\supp(\phi')| \le |\supp(\diam{\ell}\phi')| \le |\supp(\phi)| + \|\phi\|
			\]
			Similar can be shown for $\sq{\ell} \phi'$.
			% 
			\item $\fresh{x} \phi'$: By the inductive hypothesis, we have that $|\supp(\fresh{x} \phi')| + \|\fresh{x} \phi'\| \le |\supp(\phi)| + \|\phi\|$. 
			We have that $\supp(\phi') \subseteq \supp(\fresh{x} \phi')$ and $\|\phi'\| + 1= \|\fresh{x}\phi' \|$, and so:
			\[
				|\supp(\phi')| + \|\phi'\| \le |\supp(\fresh{x}\phi')| + \|\fresh{x}\phi'\|
			\]
			Similar can be shown for $\bigvee_x \hat \phi$ and $\bigwedge_x \hat \phi$.
			%
			\item $(\mu X(\vec x).\phi')(\vec u)$: By the inductive hypothesis, we have that $|\supp((\mu X(\vec x).\phi')(\vec u))| + \|(\mu X(\vec x).\phi')(\vec u)\| \le |\supp(\phi)| + \|\phi\|$. 
			We know that $\supp(\phi') \subseteq \supp(\phi)$ and $\|\phi'\| + |\vec x| = \|\phi\|$, and so:
			\[
				|\supp(\phi')| + \|\phi'\| \le |\supp(\phi)| + \|\phi\|
			\]
%			We know that $\phi' = \hat \phi\{\vec a/\vec x\}$ for some $\vec a \in \Atoms^{\ar(X)}$, so $\supp(\phi') \subseteq \supp((\mu X(\vec x).\hat \phi)(\vec u)) \cup \{\vec a\}$.
%			We also have that $\|(\mu X(\vec x).\hat \phi)(\vec u)\| = \|\phi'\| + |\vec a|$.
%			Therefore:
%			\[
%			|\supp(\phi')| \|\phi'\| \le (|\supp((\mu X(\vec x).\hat \phi)(\vec u))| + |\vec a|) + (\|(\mu X(\vec x).\hat \phi)(\vec u)\| - |\vec a|) \le |\supp(\phi)| + \|\phi\|
%			\]
			Similar can be shown for $(\nu X(\vec x).\hat \phi)(\vec u)$.
		\end{itemize}
	\end{enumerate}
\end{proof}

\subsection{Proof of \cref{thm:FnomSatPar}}

\begin{lemma}\label{lem:parityNom}
	In a game $\Gcal(\Lcal, \psi, \xi, s, H)$, if defender follows a winning strategy from some position $(s, \xi, (\mu X(\vec x).\phi)(\vec a))$,
	then there must be finitely many unfoldings of the form $(\mu X(\vec x).\phi)(\vec k)$ for $\vec k \in \Atoms^{\ar(X)}$. 
\end{lemma}
\begin{proof}
	If defender has a winning strategy from $(s, \xi, (\mu X(\vec x).\phi)(\vec a))$, then this means for every play one of the following possibilities must be true:
	\begin{enumerate}
		\item the play is finite,
		\item the play is infinite with some position with a higher even rank.
	\end{enumerate}
	For point 1, it follows that in a finite play, there will be finitely many unfoldings of $(\mu X(\vec x).\phi)(\vec a)$ as required.
	For point 2, suppose that there is an infinite play that defender wins.
	Let us say that variable $Y$ has the maximum infinitely recurring rank $d$, and as it is an even rank, $d$ is owned by defender.
	By the previous lemma, if a $(\mu X(\vec x).\phi)(\vec a)$-position reaches a $(\nu Y(\vec y).\phi')(\vec b)$-position that then reaches some $(\mu X(\vec x).\phi)(\vec k)$-position (for any $\vec k \in \Atoms^{\ar(X)}$), then $X \le_\psi Y$ and so $X$ has a higher rank than $d$.
	If this rank recurs infinitely often, then this means the play is losing for defender, and so $\mu X$ must unfold finitely many times.
\end{proof}

\begin{lemma}\label{lem:parityNom2}
	In a game $\Gcal(\Lcal, \psi, \xi, \hat s, \hat H)$, 
	if we have a $(\sigma X(\vec x).\phi')(\vec a)$-position  that reaches a $(\sigma' Y(\vec y).\phi'')(\vec b)$-position which then reaches a $(\sigma X(\vec x).\phi')(\vec c)$-position, then $X \le_\psi Y$.  
\end{lemma}%\fixednote{Update as in Lemma 81?}
\begin{proof}
	We let $\theta = \{\sigma X(\vec x).\phi'/ X\}$ and $\gamma = \{\vec a/ \vec x\}$ be substitutions and $\phi=(\sigma X(\vec x).\phi')(\vec a)$.
	In the game, the position, say, $(s, H, \xi, (\sigma X(\vec x).\phi')(\vec a))$ can make a move to $(s, H, \xi, \phi'\{\theta\}\{\gamma\})$, in particular this will include $(X(\vec c))\{\theta\}\{\gamma\} = (\sigma X(\vec x).\phi')(\vec c)$. 
	From this point, a (first) position containing $(\sigma' Y(\vec y).\phi'')(\vec b)$ can be reached.
	Then, there will be substitutions $\theta', \gamma'$ and subformula $\hat\phi''$ of $\phi$ such that $\theta \subseteq \theta'$, $\gamma \subseteq \gamma'$,
	$(\sigma'Y(\vec y).\hat \phi'')(\vec b)\{\theta'\}\{\gamma'\} = (\sigma' Y(\vec y).\phi'')(\vec b)$
	and 
	$((\sigma' Y(\vec y).\hat\phi'')(\vec b), \gamma', \theta') \in \clos(\phi)$
	(cf.~\lemmapoint{lem:freshClosProps}{2}).
	From this position, a position containing some $(\sigma X(\vec x).\phi')(\vec c)$ can be reached.
	As before,
	there must exist some $\theta' \subseteq \theta''$, $\gamma' \subseteq \gamma''$ 
	such that $\hat \phi\{\theta''\}\{\gamma''\} = (\sigma X(\vec x).\phi')(\vec c)$ and $(\hat \phi, \gamma'', \theta'') \in \clos(\phi)$. 
	Since $\hat \phi$ must be a subformula of $(\sigma X(\vec x).\phi')(\vec a)$, we can only have that $\hat \phi = X(\vec c)$, 
	meaning $X(\vec c)$ is free in $(\sigma Y(\vec y).\phi'')(\vec b)$ and so $X \le_\psi Y$.
\end{proof}

\FnomSatPar*
%	For every negation-free formula $\phi_0$,
%	$\Ucal$-variable assignment $\xi$,
%	fresh-orbit-finite F-LTS $\Lcal$ and $(s_0, H_0) \in \fstates$, then $(s_0, H_0) \in \sem{\phi_0}$ iff Defender wins from position $(s_0, H_0,\xi, \phi_0)$ in $\fgame(\Lcal, \phi_0, \xi, s_0,H_0)$.
%\end{theorem}

\begin{proof}
	Let us choose some \fhml\ formula $\phi$ with a $\Ucal$-variable assignment $\xi$ and fresh-orbit-finite F-LTS $\Lcal$
	We perform 
	induction on $\phi$ with the following base cases:
	\begin{itemize}
		\item $(s, H, \xi, a = a)$: $(s, H)\in\sem{\phi} = \fstates$, defender wins (similarly for $a \neq b$),
		\item $(s, H, \xi, a = b)$: $(s, H)\in\sem{\phi} = \emptyset$, attacker wins (similarly for $a \neq a$),
		\item $(s, H, \xi, X(\vec a))$: if $(s, H)\in \xi(X)(\vec a)$, defender wins and $(s, H) \in \sem{\phi}$, else attacker wins and $(s, H)\notin \sem{\phi}$
	\end{itemize}
	Next we perform inductive steps :
	\begin{itemize}
		\item $(s, H, \xi, \phi_1 \lor \phi_2)$:
		Let us assume $(s, H) \in \sem{\phi_1 \lor \phi_2}$.
		By the inductive hypothesis, if $(s, H) \in \sem{\phi_i}$, then defender has a winning strategy from $(s, H, \phi_i)$ for some $i \in \{1, 2\}$. 
		For any such $i$, if $s \in \sem{\phi_i}$ then $(s, H)\in \sem{\phi_1 \lor \phi_2}$, and defender has a winning strategy by the play $
		(s, H, \xi, \phi_1 \lor \phi_2) \rightarrow (s, H, \xi, \phi_i)$.
		
		Next, let us assume that defender has a winning strategy from $(s, H, \xi, \phi_1 \lor \phi_2)$. 
		By the inductive hypothesis, if defender wins from $(s, H, \xi, \phi_i)$ then $(s, H)\in \sem{\phi_i}$ for some $i \in \{1, 2\}$.
		As defender wins from $(s, H, \xi, \phi_1 \lor \phi_2)$, defender wins from $(s, H, \xi, \phi_i)$ for some $i \in \{1, 2\}$, then $(s, H) \in \sem{\phi_i} \subseteq \sem{\phi_1 \lor \phi_2}$ as required.
		
		The proof for $\phi_1 \land \phi_2$ is similar.
		%
		\item $(s, \xi, \bigvee_x \phi')$: 
		By the inductive hypothesis, if $(s, H)\in \sem{\phi'\{a/x\}}$, then defender has a winning strategy from $(s, H, \xi, \phi\{a/x\})$ for some $a \in \Atoms$.
		For any such $a$, if $(s, H)\in \sem{\phi'\{a/x\}}$ then $(s, H) \in \sem{\bigvee_x \phi'}$, and defender has a winning strategy by the move $(s, H, \xi, \bigvee_x \phi') \rightarrow (s, H, \xi, \phi'\{a/x\})$.
		
		Next, let us assume that defender has a winning strategy from $(s, H, \xi, \bigvee_x \phi')$. 
		By the inductive hypothesis, if defender wins from $(s, H, \xi, \phi'\{a/x\})$ then $(s, H) \in \sem{\phi'\{a/x\}}$ for some $a \in \Atoms$.
		As defender wins from $(s, H, \xi, \bigvee_x \phi')$, defender wins from $(s, H, \xi, \phi\{a/x\})$ for some $a \in \Atoms$, then $(s, H) \in \sem{\phi\{a/x\}} \subseteq \sem{\bigvee_x \phi'}$ as required.
		
		The proof for $\bigwedge_x \phi'$ and $\fresh{x} \phi'$ is similar.
		%
		\item $(s, H,\diam{\ell} \phi')$: If $(s, H)\in \sem{\diam{\ell} \phi'}$, then there exists some $(s', H')$ such that $(s, H) \xrightarrow{\ell} (s', H')$ and $(s', H')\in \sem{\phi'}$. 
		By the inductive hypothesis, if $(s', H')\in \sem{\phi'}$ then defender has a winning strategy from $(s', H', \xi, \phi')$.
		As this position is played by defender, it suffices to choose $(s', H', \xi, \phi')$, meaning defender has a winning strategy.  
		
		Next, let us assume that defender has a winning strategy from $(s, H, \xi, \diam{\ell} \phi')$. 
		This means there is some move $(s, H, \xi, \diam{\ell} \phi') \to (s', H', \xi, \phi')$ with $(s, H) \xrightarrow{\ell} (s', H')$. 
		By the inductive hypothesis, if defender wins from $(s', H', \xi, \phi')$ then $(s', H') \in \sem{\phi'}$, thus:
		\[
			(s, H) \in \{(s_1, H_1) \in \fstates \mid \exists (s_1, H_1) \xrightarrow{\ell} (s_2, H_2).(s_2, H_2) \in \sem{\phi'}\} = \sem{\phi}
		\]	
%		which implies that $(s, H) \in \sem{\diam{\ell}\phi'}$ as required.
% \hbnote{@NT: changed this as I believe this makes more sense, please check if you can!}
%		By the inductive hypothesis, if defender wins from $(s, H, \xi, \phi')$ then $(s, H) \in \sem{\phi'}$.
%		As defender wins from $(s, H, \xi, \diam{\ell} \phi')$, defender wins from $(s, H, \xi, \phi')$, then $(s, H) \in \sem{\phi'} \subseteq \sem{\diam{\ell}\phi'}$ as required.
%		
		The proof for $\sq{\ell} \phi'$ is similar.
		
	\end{itemize}
	%
	The more interesting cases are the fixpoint operators.
	Starting with $\mu$, we want to show that, for any $(s', H')$:
	\[
	(s', H')\in \sem{(\mu X(\vec x).\phi')(\vec a)} \text{ iff defender wins from } (s', H', \xi, (\mu X(\vec x).\phi')(\vec a)) \text{ in } \Gcal(\Lcal, \phi, \xi, s, H)
	\]
	We begin with left-to-right implication, assuming that $(s', H') \in \sem{(\mu X(\vec x).\phi')(\vec a)}$. 
	We define the following function:%\fixednote{typo: $\phi'$?}
	\[
	F = \lambda f.\lambda \vec b.\sem[{\xi[X \mapsto f]}]{\phi'\{\vec b / \vec x\}}
	\]
	We denote that for any $k$ that $\xi^k = \xi[X \mapsto F^k(\emptyset)]$.
	Furthermore, we denote $\Gcal_{k}$ to be the game $\Gcal(\Lcal, \phi, \xi^k, s, H)$ for all $k$, and denote simply $\Gcal$ for $\Gcal(\Lcal, \phi, \xi, s, H)$.
	By definition, we have that:
	\[
	\sem{(\mu X(\vec x).\phi')(\vec a)} = \lfp(F)(\vec a) = (\bigcup_{k \in \omega} F^k(\emptyset))(\vec a)
	\]
	We show that for all $k$ and $\vec c$, if defender wins from position $(s', H', \xi^k, \phi'\{\vec c/\vec x\})$ in $\Gcal_k$,
	then defender wins from $(s', \xi, (\mu X(\vec x).\phi')(\vec c))$.
	Note that the latter has unique next position $(s', H', \xi, (\phi'\{\mu X(\vec x).\phi' / X\}\{\vec c / \vec x\}))$. 
	We perform induction on $k$. 
	If $k = 0$, then moves from $(s', \xi^0, \phi'\{\vec c / \vec x\})$ in $\Gcal_0$ never reach a position $(s'',H'', \xi^0, X(\vec d))$ for some $\vec d$. 
	This is because $F^0(\emptyset)(\vec d) = \emptyset$, and so it is not possible to win from $(s'', H'', \xi^0, X(\vec d))$.
	Thus, the same strategy used to win from $(s', H', \xi^0, \phi'\{\vec c / \vec x\})$ can be used to win from $(s', H', \xi, (\mu X(\vec x).\phi')(\vec c))$ in $\Gcal$.
	Suppose that defender wins from $(s', H', \xi^{k + 1}, \phi'\{\vec c / \vec x\})$ in $\Gcal_{k + 1}$, then we obtain the following strategy for $(s', H', \xi, (\phi'\{\mu X(\vec x).\phi' / X\})\{\vec c / \vec x\})$:
	play as normal in $\Gcal$ from $(s', H', \xi, (\phi'\{\mu X(\vec x).\phi' / X\})\{\vec c / \vec x\})$.
	If some $(s'', H'', \xi, (\mu X(\vec x).\phi')(\vec d))$ is reached, then $(s'', H'', \xi^{k + 1}, X(\vec d))$ is reached in $\Gcal_{k + 1}$ from $(s', H', \xi^{k+1}, \phi'\{\vec c / \vec x\})$. 
	By definition, defender is winning from $(s'', H'', \xi^{k + 1}, X(\vec d))$ in $\Gcal_{k+1}$ iff $(s'', H'')\in \sem[\xi^{k + 1}]{X(\vec d)}$.
	Then $(s'', H'') \in \xi^{k + 1}(X)(\vec d) = (F^{k + 1}(\emptyset))(\vec d) = \sem[\xi^k]{\phi'\{\vec d / \vec x\}}$.
	Hence, using the outer induction hypothesis we deduce that defender wins from $(s'', H'',\xi^k, \phi'\{\vec d / \vec x\})$ in $\Gcal_k$, and so by the inner inductive hypothesis, defender has a winning strategy from $(s'', H'', \xi, (\mu X(\vec x).\phi')(\vec d))$ in $\Gcal$, thus this completes the proof of the inner claim.
	\\
	As $(s', H') \in \sem{(\mu X(\vec x).\phi')(\vec a)}$, then $(s', H') \in \bigcup_{j \in \omega} (F^j(\emptyset))(\vec a)$ so $(s', H') \in (F^k(\emptyset))(\vec a)$ for some $k \in \omega$. 
	Therefore, we have that $k > 0$ and $(s', H') \in (F(F^{k - 1}(\emptyset)))(\vec a) = \sem[\xi^k]{\phi'\{\vec a / \vec x\}}$.
	Thus, using the outer inductive hypothesis, defender wins from $(s, H, \xi^k, \phi'\{\vec a / \vec x\})$ in $\Gcal_k$ and so defender wins from $(s', H', \xi, (\mu X(\vec x).\phi')(\vec a))$ in $\Gcal$.
	
	Next, suppose defender wins from $(s', H', \xi, (\mu X(\vec x).\phi')(\vec a))$ in $\Gcal$. 
	By Lemma~\ref{lem:parityNom}, and as defender is winning, every play in this game must have finitely many unfoldings of the form $(\mu X(\vec x).\phi')(\vec c)$ for $\vec c \in \Atoms^{\ar(X)}$.
	Let us consider the winning strategy for defender from $(s', H', \xi, (\mu X(\vec x).\phi')(\vec a))$.
	This can be interpreted as a sub-graph of the full game where, in defender positions, there is a unique successor position respecting the winning strategy (in attacker positions, all successor positions are present). 
	We call positions of the form $(s'', H'', \xi, (\mu X(\vec x).\phi')(\vec c))$ \emph{critical}. 
	By the previous argument, each critical position in the graph, in any play, can each reach finitely many other critical positions. 
	We define the height of every critical position $c$ (denoted $h(c)$) as the maximum number of critical positions reached during any play from $c$. 
	We show that for each $c = (s'', H'', \xi, (\mu X(\vec x).\phi')(\vec c))$, then $(s'', H'')\in \sem[\xi^{h-1}]{\phi'\{\vec c / \vec x\}}$ with $h = h(c)$. 
	This then implies that $(s', H') \in \sem[\xi^{h - 1}]{\{\phi'\{\vec c / \vec x\}} \subseteq \sem{(\mu X(\vec x).\phi')(\vec a)}$.
	We perform induction on $h$ with $h \ge 1$ (a critical position can always reach itself). 
	For $h = 1$, the defender wins from $(s'',H'', \xi, (\mu X(\vec x).\phi')(\vec c))$ without reaching any other critical position, therefore defender also wins from $(s'', \xi^1, \phi'\{\vec c / \vec x\})$ in $\Gcal_1$.
	Thus by the outer inductive hypothesis, $(s'', H'') \in \sem{\phi'\{\vec c / \vec x\}}$.
	Suppose that $c$ has height $h + 1$. 
	Then the defender wins from $(s'', H'', \xi^h, \phi\{\vec c / \vec x\})$ in $\Gcal_h$ as they can play as for $(s'',H'', \xi, (\mu X(\vec x).\phi')(\vec c))$ in $\Gcal$ and if they reach some $(s''', H''', \xi^h, X(\vec d))$ then, in the game with a winning strategy $\Gcal$, this reaches the winning position $(s''', H''', \xi, (\mu X(\vec x).\phi')(\vec d))$ with height less than or equal to $h$. 
	By the inner inductive hypothesis, we have that $(s''', H''')\in \sem[\xi^{h-1}]{\phi'\{\vec d / \vec x\}} = (F^h(\emptyset))(\vec d) = \xi^h(X)(\vec d)$, thus defender wins from $(s''', \xi^h, X(\vec d))$ in $\Gcal_h$ and so defender wins from $(s', H', \xi^h, \phi'\{\vec c / \vec x\})$. 
	By applying the outer inductive hypothesis we obtain $(s'', H'')\in \sem[\xi^h]{\phi'\{\vec c / \vec x\}}$.  
	%
	%
	%
	
	Next, we examine the case of $\nu$. 
	Similarly to $\mu$, we want to show that:
	% for any $\vec a$:
	\[
	(s', H')\in \sem{(\nu X(\vec x).\phi')(\vec a)} \text{ iff defender wins from }(s', H', \xi, (\nu X(\vec x).\phi')(\vec a))\text{ in }\Gcal(\Lcal, \phi, \xi, s, H)
	\]
	The latter position has the following unique successor position
	$(s', H', \xi, (\phi'\{\nu X(\vec x).\phi'/X\})\{\vec a/\vec x\})$.
	We begin with left-to-right implication, assuming that $(s', H')\in \sem{(\nu X(\vec x).\phi')(\vec a)}$.
	We note that:
	\begin{align*}
		&\sem{(\nu X(\vec x).\phi')(\vec a)} = \sem[\xi']{\phi'\{\vec a / \vec x\}}\\
		&\text{ with } \xi' = \xi[X \mapsto \sem{(\nu X(\vec x).\phi')}]
	\end{align*}
	By the inductive hypothesis, $(s', H') \in \sem[\xi']{\phi'\{\vec a / \vec x\}}$ implies that defender wins from $(s', H', \xi', \phi'\{\vec a / \vec x\})$.
	We build a winning strategy 
	from $(s', H',\xi,(\nu X(\vec x).\phi')(\vec a))$
	for the defender as follows: the defender follows the strategy for $(s', H', \xi', \phi'\{\vec a / \vec x\})$.
	If some position $(s'', H'', \xi, (\nu X(\vec x).\phi')(\vec b))$ is reached, this would correspond to reaching $(s'', H'', \xi', X(\vec b))$ in the game the defender wins. 
	As the defender is winning, this means $(s'', H'')\in \xi'(X)(\vec b) = \sem{(\nu X(\vec x).\phi')(\vec b)}$, hence $(s'', H'')\in \sem[\xi']{\phi'\{\vec b / \vec x\}} = \sem{(\nu X(\vec x).\phi')(\vec b)}$.
	Then, by the inductive hypothesis, defender wins from $(s'', H'', \xi, \phi'\{\vec b / \vec x\})$. 
	We continue in our game to $(s'', H'',\xi, (\phi'\{\nu X(\vec x).\phi'/X\})\{\vec b / \vec x\})$ following the winning strategy for $(s'', H'', \xi, \phi\{\vec b / \vec x\})$, and so on.
	This strategy is winning as:
	\begin{enumerate}
		\item if there are finitely many unfoldings of the form $X(\vec k)$, then the last $(s'', H'', \xi, \phi\{\vec k / \vec x\})$ is winning for the defender,
		\item otherwise, it must be that 
		some $X(\vec k)$ has infinitely many unfoldings, so defender wins as $\Omega((\nu X(\vec x).\phi')(\vec k))$ is even and must be the highest infinitely repeating rank (c.f Lemma~\ref{lem:parityNom2}).
	\end{enumerate}
	%
	Lastly, we assume defender has a winning strategy from $(s', H', \xi, (\nu X(\vec x).\phi')(\vec a))$. 
	We redefine for any $j$ that $\xi^j = \xi[X \mapsto F^j(\fstates)]$.
%	As before, we denote $\Gcal_{k}$ to be the game $\Gcal(\Lcal, \phi, \xi^k, s, H)$ for all $k$, and denote simply $\Gcal$ for $\Gcal(\Lcal, \phi, \xi, s, H)$. 
	By definition we have that:
	\[
		\sem{(\nu X(\vec x).\phi')(\vec a)} = \gfp(F)(\vec a) = (\bigcap_{j \in \omega} F^j(\Scal))(\vec a)
	\]
	We claim that the defender wins from $(s', H', \xi^j,\phi\{\vec a / \vec x\})$ for all $j$. 
	For $j = 0$, we just play as in $(s', H', \xi, (\phi\{\nu X(\vec x).\phi' / X\})\{\vec a / \vec x\})$, and if we reach some $(s'', H', \xi^0, X(\vec c))$, then defender wins. 
	This is because $(s'', H'') \in \xi^0(X)(\vec c) = (F^0(\Scal))(\vec c) = \Scal$. 
	If some $(s'', H'', \xi^{j + 1}, \phi\{\vec c / \vec x\})$ is reached, then we need to show that $(s'', H'')\in \xi^{j+1}(X)(\vec c) = (F^{j + 1}(\fstates))(\vec c)$. 
	In the game from $(s', H', \xi, (\phi'\{\nu X(\vec x) / X\})\{\vec a / \vec x\})$, we would reach $(s'', H'', \xi, (\nu X(\vec x).\phi')(\vec c))$. 
	By the inductive hypothesis, defender wins from $(s'', H'',\xi, \phi\{\vec c / \vec x\})$, and by the outer inductive hypothesis, $(s'', H'')\in \sem[\xi^j]{\phi\{\vec c / \vec x\}} = \xi^{j + 1}(X)(\vec c)$ as required.
	Thus, by the outer inductive hypothesis, $(s', H')\in \sem[\xi^j]{\phi\{\vec c / \vec x\}}$ for all j, and therefore 
	\[
	(s', H')\in (\bigcap_{j \in \omega}F^j(\Scal))(\vec a) = \sem{(\nu X(\vec x).\phi')(\vec a)}
	\]
\end{proof}

%\beyondthisline
%\begin{proof}
%    The proof is symmetric to Theorem~\ref{thm:nomSatPar}. We show the inductive step for the additional $\fresh{x}$ connective:
%    \begin{itemize}
%        \item $(s, H, \xi, \fresh{x}\phi')$: 
%        By the inductive hypothesis, if $(s, H)\in \sem{\phi'\{a/x\}}$, then defender has a winning strategy from $(s, H, \xi, \phi\{a/x\})$ for some $(a, H) \in \Atoms \setminus H$.
%        For any such $a$, if $(s, H)\in \sem{\phi'\{a/x\}}$ then $(s, H) \in \sem{\fresh{x} \phi'}$, and defender has a winning strategy by the move $(s, H, \xi, \fresh{x} \phi') \rightarrow (s, H, \xi, \phi'\{a/x\})$.
%
%        Next, let us assume that defender has a winning strategy from $(s, H, \xi, \fresh{x} \phi')$. 
%        By the inductive hypothesis, if defender wins from $(s, H, \xi, \phi'\{a/x\})$ then $(s, H) \in \sem{\phi'\{a/x\}}$ for some $a \in \Atoms \setminus H$.
%        As defender wins from $(s, H, \xi, \fresh{X} \phi')$, defender also wins from $(s, H, \xi, \phi\{a/x\})$ for some $a \in \Atoms\setminus H$, then $(s, H) \in \sem{\phi\{a/x\}} \subseteq \sem{\fresh{x} \phi'}$ as required.
%    \end{itemize}
%  \end{proof}

\subsection{Proof of \cref{lem:boundingisnombisim}}

\boundingisnombisim
\begin{proof}
  Let us write $\Rcal$ for $\Rcal_{(\Lcal, \phi_0, s_0, H_0)}$.
Let us choose some $((s, H_1, \phi),(s, H_2, \phi)) \in \Rcal_{(\Lcal, \phi_0, s_0, H_0)}$. By definition, the rank and polarity of $(s, H_1, \phi)$ and $(s, H_2, \phi)$ are the same.
    %We how that each move from $(s, H_1, \phi) \rightarrow (s', H_1', \phi')$ has a matching move $(s, H_2, \phi) \rightarrow (s', H_2', \phi')$ such that $((s', H_1', \phi'),(s', H_2', \phi')) \in \Rcal_{(\Lcal, \hat \phi, \hat s, \hat H)}$. 
    %For every $(P_1, P_2) \in \Rcal_{(\Lcal, \hat \phi, \hat s, \hat H)}$, as $P_1$ has the same formula and rank as $P_2$, 
Let us introduce the following notation for economy (where $ H \subseteq \Atoms$) :\
        $ H_{(-b+a)} = ( H \setminus \{b\}) \cup \{a\}$.\\
        We next look at the possible moves.
We perform a case analysis on $\phi$:
    \begin{itemize}
        \item $u = v$: neither $(s, H_1, \phi)$ nor $(s, H_2, \phi)$ can make a move.
        %, therefore, $\rel{(s, H_1, \phi)}{(s, H_2, \phi)}$.
        %
        \item $\phi_1 \lor \phi_2$: we have that $(s, H_1, \phi) \rightarrow (s, H_1, \phi_i)$ and $(s, H_2, \phi) \rightarrow (s, H_2, \phi_i)$ for $i \in \{1, 2\}$. 
        For any such $i$, as the history does not change, if $|H_2| \le \S $ then $((s, H_1, \phi_i), (s, H_2, \phi_i)) \in \Rcal$ for $i \in \{1,2\}$. 
        Now suppose $|H_2| = \S + 1$. 
        As $\supp(\phi_i) \subseteq \supp(\phi)$, it follows that $\supp(\phi_i, s) \cap H_1 \subseteq \supp(\phi, s) \cap H_1 \subseteq H_2 \subseteq H_1$, meaning $((s, H_1, \phi_i),(s, H_2, \phi_i)) \in \Rcal$ for $i \in \{1,2\}$ as required.
        Similar can be shown for $\phi_1 \land \phi_2$.
        %
        \item $\fresh{x} \phi'$ the possible moves in $\fgame$ are $(s, H_1, \phi)\rightarrow (s, H_1, \phi'\{a/x\})$ for all $a \# \phi,H_1$, hence each such move can be matched by some $(s, H_2, \phi) \rightarrow (s, H_2, \phi'\{a/x\})$ and we have that $a \notin H_2 \subseteq H_1$ hence $((s, H_1, \phi'\{a/x\}),(s, H_2, \phi'\{a/x\})) \in \Rcal$. 
        \\
        Conversely, the possible moves in $\Gcal_\S$ are $(s, H_2, \phi) \rightarrow (s, H_2, \phi'\{a/x\})$ for $a \# \phi,H_2$. 
        If $a \notin H_1$, then as before this can be matched by $(s, H_1, \phi)\rightarrow (s, H_1, \phi'\{a/x\})$ with $((s, H_1, \phi'\{a/x\})$, $(s, H_2, \phi'\{a/x\})) \in \Rcal$.
        Otherwise, {and this is the more interesting case}, suppose $a \in H_1 \setminus H_2\cup\supp(\phi)$.
Then, we can select some $b \notin H_1$, meaning $b \notin H_2$ and $(a\ b)\cdot (s, H_2, \phi'\{a/x\}) = (s, H_2, \phi'\{b/x\})$. This can be matched by $(s, H_1, \phi) \rightarrow (s, H_1, \phi'\{b/x\})$ with $((s, H_1, \phi'\{b/x\}), (s, H_2, \phi'\{b/x\})) \in \Rcal$ as required. 
        %
        \item $\bigvee_x \phi'$: 
        if $|H_2| \le \S$, then $(s, H_1, \bigvee_x \phi') \rightarrow (s, H_1, \phi'\{a/x\})$ can be matched by the move $(s, H_2, \bigvee_x \phi') \rightarrow (s, H_2, \phi'\{a/x\})$ with $((s, H_1, \phi'\{a/x\}), (s, H_2, \phi'\{a/x\})) \in \Rcal$ as required. And similarly for each $(s, H_2, \bigvee_x \phi') \rightarrow (s, H_2, \phi'\{a/x\})$.
        \\
        Suppose $|H_2| = \S+1$, and
        let $(s, H_1, \bigvee_x \phi') \rightarrow (s, H_1, \phi'\{a/x\})$ for some $a \in \Atoms$.
Note that $\supp(\phi', s) \cap H_1 \subseteq H_2 \subseteq H_1$. We do a case analysis on $a$.
        \begin{itemize}
            \item 
        If $a \in H_2$ or $a\notin H_1$, then this means $a$ is in both histories, or neither, so it suffices to again match with the move $(s, H_2, \bigvee_x \phi') \rightarrow (s, H_2, \phi'\{a/x\})$ with $((s, H_1, \phi'\{a/x\})$, $(s, H_2, \phi'\{a/x\})) \in \Rcal$. 
        \item
        If $a \in H_1 \setminus H_2$, then we no longer have 
        $((s, H_1, \phi'\{a/x\}),(s, H_2, \phi'\{a/x\})) \in \Rcal$. 
        Note that then $a\notin\supp(s,\phi')$ by definition of $\Rcal$. 
        We pick $b\in H_2\setminus\supp(s,\phi')$ and make the move
        $(s, H_2, \bigvee_x \phi') \rightarrow (s, H_2, \phi'\{b/x\})$. We have $(a\ b)\cdot(s, H_2, \phi'\{b/x\})=(s, H_{2(-b+a)}, \phi'\{a/x\})$ and
        $((s, H_1, \phi'\{a/x\}),(s, H_{2(-b+a)}, \phi'\{a/x\})) \in \Rcal$, as required.
                \end{itemize}
        Conversely, suppose $(s, H_2, \bigvee_x \phi')\to(s, H_2, \phi'\{a/x\})$. We do case analysis on $a$.
        \begin{itemize}
        \item If $a \in H_2$ or $a \notin H_1$, then
          $(s, H_1, \bigvee_x \phi')$ $ \rightarrow (s, H_1, \phi'\{a/x\})\Rcal(s, H_2, \phi'\{a/x\})$.
            %
            \item If $a \in H_1 \setminus H_2$, then let us choose the transition $(s, H_1, \bigvee_x \phi') \to (s, H_1, \phi'\{b/x\})$ for some $b \notin H_1 \cup \supp(\phi', s)$.
            Then, we have that $(s, H_1, \phi'\{b/x\}) \nomeq (s, H_{1(-a+b)}, \phi'\{a/x\})$ and $((s, H_{1(-a+b)}, \phi'\{a/x\}),(s, H_2, \phi'\{a/x\})) \in \Rcal$ as required.
        \end{itemize}
        % Similar can be shown for $\bigwedge_x \phi'$.
        %
        \item $\diam{t, \vec a} \phi'$: in $\fgame$, the possible moves are $(s, H_1, \phi) \rightarrow (s', H_1', \phi')$ for all $(s', H_1')$ such that $(s, H_1) \xrightarrow{t, \vec a} (s', H_1')$.
        By definition, we know $\supp(s') \subseteq \supp(s) \cup \{\vec a\}$ and $H_1' = H_1 \cup \{\vec a\}$. 
        If $|H_1'| \le \S$, then this is matched by some move $(s,H_2,\phi) \rightarrow (s', H_2', \phi')$ with $H_2' = H_2 \cup \{\vec a\}$ and $((s', H_1', \phi'),(s',H_1',\phi')) \in \Rcal$. 
        Otherwise, we can see that $|H_2 \cup \{\vec a\}| > \S$. This can be matched by some move $(s,H_2,\phi) \rightarrow (s', H_2', \phi')$ with $H_1' \cap \supp(\phi', s') \subseteq H_2' \subseteq H_1'$ and $|H_2'| = \S + 1$. 
        As $|\supp(\phi',s')| \le \S$, we have that such a $H_2'$ exists and $((s', H_1', \phi'),(s',H_2',\phi')) \in \Rcal$ as required.
        \\
        Conversely, let $(s, H_2, \phi) \rightarrow (s', H_2', \phi')$ with $s\xrightarrow{t, \vec a} s'$. One of the following is the case:
        \begin{itemize}
            \item $|H_2 \cup \{\vec a\}| \le \S$, then $H_2' = H_1 \cup \{\vec a\}$. 
            This can be matched by the move $(s, H_1, \phi) \rightarrow (s', H_1', \phi')$ with $H_1' = H_1 \cup \{\vec a\}$ and $((s', H_1', \phi'),(s',H_2',\phi')) \in \Rcal$.
            \item $|H_2 \cup \{\vec a\}| > \S$ and $(H_2 \cup \{\vec a\}) \cap \supp(s', \phi') \subseteq H_2' \subseteq H_2 \cup \{\vec a\}$ with $|H_2'| = \S + 1$. 
            As $\supp(s') \subseteq \supp(s) \cup \{\vec a\}$ and $\supp(\phi) = \supp(\phi') \cup \{\vec a\}$, we have:
            \begin{align*}
                (H_1 \cup \{\vec a\}) \cap \supp(s', \phi') &= ((H_1\setminus\{\vec a\}) \cap \supp(s', \phi')) \cup (\{\vec a\} \cap \supp(s', \phi'))
                \\
                &\subseteq (H_2 \cap \supp(s', \phi')) \cup (\{\vec a\} \cap \supp(s', \phi'))\subseteq H_2' 
                % \\&= (H_2 \cup \{\vec a\}) \cap \supp(s', \phi') \subseteq H_2' 
            \end{align*}
            Thus, the move can be matched by $(s, H_1, \phi) \rightarrow (s', H_1', \phi')$ with $H_1' = H_1 \cup \{\vec a\}$ and $((s', H_1', \phi'),(s',H_2',\phi')) \in \Rcal$.
        \end{itemize}
% Similar can be shown for $\sq{t, \vec a} \phi'$.
        %
        \item $(\mu X(\vec x).\phi')(\vec a)$: in $\fgame$, the only move is $(s, H_1, \phi) \rightarrow (s, H_1, (\phi'\{\mu X(\vec x).\phi'/X\})\{\vec a/\vec x\})$. As $H_1$ does not change and $\supp(\phi, s) = \supp((\phi'\{\mu X(\vec x).\phi'/X\})\{\vec a/\vec x\})$, then this can be matched by the only move in $\Gcal_\S$, i.e.\ $(s, H_2, \phi) \rightarrow (s, H_2, (\phi'\{\mu X(\vec x).\phi'/X\})\{\vec a/\vec x\})$ with $((s, H_1, (\phi'\{\mu X(\vec x).\phi'/X\})\{\vec a/\vec x\}), (s, H_2, (\phi'\{\mu X(\vec x).\phi'/X\})\{\vec a/\vec x\})) \in \Rcal$. 
        % Similar can be shown for $(\nu X(\vec x).\phi')(\vec a)$.
    \end{itemize} 
The remaining cases can be shown similarly to the ones above.
  \end{proof}

\subsection{Proof of \cref{lem:nombisimequiv}}
\nombisimequiv*
\begin{proof}
  By symmetry, it suffices to show the left-to-right direction. Suppose $(P_1,P_2)\in R$ for some bisimulation $R$ and let $\Theta_1$ be a winning strategy for Defender from $P_1$. We assume the existence of a choice function $\iota$ such that, for each $(P_1',P_2')\in R$:
  \begin{itemize}
\item    if $P_1',P_2'$ are Attacker positions then, for each move $P_2'\to P_2''$, $\iota(P_1',P_2',P_2'')=P_1''$ such that $P_1'\to P_1''$ and $(P_1'',P_2'')\in R$; 
\item    if $P_1',P_2'$ are Defender positions then, for each move $P_1'\to P_1''$, $\iota(P_1',P_2',P_1'')=P_2''$ such that $P_2'\to P_2''$ and $(P_1'',P_2'')\in R$.
    \end{itemize}
While our strategies are positional (i.e.\ in $V_D\rightharpoonup E$), we next define a strategy $\Theta_2:V^*\times V_D\to E$, i.e.\ one that looks at the whole play. It is known that if there is such a winning strategy for Defender from a root position then there is also a positional one~\cite{EmersonJutla91}. Let us call a play $\Pi_2=P_{2,0},P_{2,1},\dots,P_{2,n}$ in $\Gcal_2$ with $P_{2,0}=P_2$ to be \emph{$(\Theta_1,R)$-compatible} if there is a unique play 
$\Pi_1=P_{1,0},P_{1,1},\dots,P_{1,n}$ in $\Gcal_1$, called \emph{the $\Gcal_1$-match of $\Pi_2$}, such that $P_{1,0}=P_1$ and, for each $i<n$, $(P_{1,i+1},P_{2,i+1})\in R$ and:
\begin{itemize}
\item if $P_{1,i},P_{2,i}$ are Attacker positions then $P_{1,i+1}=\iota(P_{1,i},P_{2,i},P_{2,i+1})$;
\item if $P_{1,i},P_{2,i}$ are Defender positions then $\Theta_1(P_{1,i})=P_{1,i}\to P_{1,i+1}$ and $P_{2,i+1}=\iota(P_{1,i},P_{2,i},P_{1,i+1})$.
\end{itemize}
We let the strategy $\Theta_2$ be defined on \emph{$(\Theta_1,R)$-compatible} plays ending in Defender positions, by setting
$\Theta_2(P_{2,0},P_{2,1},\dots,P_{2,n})=\iota(P_{1,n},P_{2,n},\Theta_1(P_{1,n}))$, where $P_{1,0},P_{1,1},\dots,P_{1,n}$ the $\Gcal_1$-match of $P_{2,0},P_{2,1},\dots,P_{2,n}$. Note that, for all $i$, $(P_{1,i},P_{2,i})\in R$.
\\
Let now $\Pi_2=P_{2,0},P_{2,1},\dots$ be a complete or infinite play in $\Gcal_2$ starting from $P_2$ and following $\Theta_2$ for Defender. We claim that every finite prefix $P_{2,0},P_{2,1},\dots,P_{2,n}$ is $(\Theta_1,R)$-compatible. Indeed, if $P_{2,0},P_{2,1},\dots,P_{2,i}$ is $(\Theta_1,R)$-compatible with $\Gcal_1$-match $P_{1,0},P_{1,1},\dots,P_{1,i}$ then, depending on the polarity of $P_{1,i},P_{2,i}$, we get a unique  $P_{1,0},P_{1,1},\dots,P_{1,i+1}$ such that $P_{2,0},P_{2,1},\dots,P_{2,i+1}$ is $(\Theta_1,R)$-compatible. Hence, if $\Pi_2$ is complete then Defender wins as the corresponding play $\Pi_1$ must be winning by hypothesis. If, on the other hand, $\Pi_2$ is infinite then we can define the infinite play $\Pi_1$ by setting, for each $n$, $P_{1,0},P_{1,1},\dots,P_{1,n}$ to be the $\Gcal_1$-match of $P_{2,0},P_{2,1},\dots,P_{2,n}$. As Defender must win $\Pi_1$, they also win $\Pi_2$.
\end{proof}

\subsection{Proof of \cref{lem:orbitGameSize}}\label{app:orbitGameSize}



% \begin{definition}
% 	Given a fresh-orbit-finite F-LTS $\Lcal=\langle \Scal, L, \to \rangle$ with a set of configurations $\fstates$, we define its \emph{non-deterministic branching} to be as follows:
% %	\ntnote{\scriptsize we need FLTS; also, labels can be more general that single a's}
% 	\[
% 	\ndbb(\Lcal) = \max\{
% 	\ |\{(s', H') \mid (s, H) \xrightarrow{\ell} (s', H')\}|\  
% 	\mid (s, H) \in \fstates, \ell \in L\
% 	\}
% 	\]
% 	that is, the maximum number of ways any name can be taken from any configuration.
% %	Furthermore, we define $\max(\Lcal)$ to be the maximum size of any configuration in $\Lcal$, that is, $\max\{|s| + |H| \mid (s, H) \in \fstates\}$.
% %	\ntnote{\scriptsize as per previous comment, there may be no need for max(L)}
% \end{definition}

%\begin{lemma}\label{lem:cartesianOrbitSize}
%	Let $A, B$ be nominal sets. 
%	Then, $|\orb(A \times B)| \le |\orb(A)|\cdot|\orb(B)|\cdot {n_1}^{n_2}(1 + o(1))$ where
%	$n_1 = \max(\regindex(A), \regindex(B))$, $n_2 = \min(\regindex(A),\regindex(B))$
%	and $o(1)$ is an asymptotically small constant. 
%\end{lemma}
%
%\begin{proof}
%	Let $x_{O}$ range over orbit representatives in $\orb(A), \orb(B)$ such that for any $O_A \in \orb(A)$, $x_{O_A} \in O_A$, and similarly for $B$.
%	We define the following function:
%	\[
%	F(O_A, O_B) = \bigcup_{\pi' \in \Perm}\{\pi \cdot (x_{O_A}, \pi'\cdot x_{O_B}) \mid \pi \in \Perm\}
%	% \\ &= \bigcup_{\pi' \in \fix(\Atoms \setminus \supp(X_A, X_B))}\{\pi \cdot (X_A, \pi'\cdot X_B) \mid \pi \in \Perm\}
%	\]
%	Where $O_A \in \orb(A)$ and $O_B \in \orb(B)$. 
%	Without loss of generality, let us assume that $\supp(x_{O_A}) \cap \supp(x_{O_B}) = \emptyset$. 
%	Then, we can say that $\{\pi \circ \pi' \mid \pi, \pi' \in \Perm\} = \{\pi \circ \pi' \mid \pi \in \Perm, \pi' \in \fix(\Atoms \setminus \supp(x_{O_A}, x_{O_B})\}$ as both sets are closed under permutation.  
%	Let us fix $\vec b_0$ to be some sequencing of $\supp(x_{O_B})$ and let us write $\vec b' \lesssim \vec b_0$ to indicate $\vec b'$ is a sub-sequence of $\vec b_0$.
%	From this, we can say:
%	\begin{align*}
%		F(O_A, O_B) &= \bigcup_{\pi' \in \fix(\Atoms \setminus \supp(x_{O_A}, x_{O_B}))}\{\pi \cdot (x_{O_A}, \pi'\cdot x_{O_B}) \mid \pi \in \Perm\} \\
%		&= \bigcup_{\substack{\vec b_0 \lesssim \vec b,\\ \{\vec a\} \subseteq \supp(x_{O_A})}} \{\pi \cdot (x_{O_A}, \move{\vec b}{\vec a}\cdot x_{O_B}) \mid \pi \in \Perm\} \\
%		&= \bigcup_{\substack{\vec b_0 \lesssim \vec b,\\ \{\vec a\} \subseteq \supp(x_{O_A})}} \orbs((x_{O_A}, \move{\vec b}{\vec a}\cdot x_{O_B}))
%	\end{align*}
%	Next we can write $\orb(A \times B)$ in terms of $F$ as follows:
%	\begin{align*}
%		\orb(A \times B) &= \{\{(\pi \cdot x_{O_A},\pi'\cdot x_{O_B}) \mid \pi, \pi' \in \Perm\} \mid O_A \in \orb(A), O_B \in \orb(B)\}\\
%		&= \{\{\pi\cdot (x_{O_A},\pi^{-1}\cdot \pi'\cdot x_{O_B}) \mid \pi, \pi' \in \Perm\} \mid O_A \in \orb(A), O_B \in \orb(B)\}\\
%		&= \{\{\pi\cdot (x_{O_A},\pi''\cdot x_{O_B}) \mid \pi, \pi'' \in \Perm\} \mid O_A \in \orb(A), O_B \in \orb(B)\}\\
%		&= \bigcup_{\substack{O_A \in \orb(A),\\ O_B \in \orb(B)}}F(O_A, O_B)\\
%		\therefore\ |\orb(A \times B)| &\le \sum_{\substack{O_A \in \orb(A),\\ O_B \in \orb(B)}}|F(O_A, O_B)| 
%	\end{align*}
%	It remains to calculate $|F(O_A, O_B)|$ for any $O_A \in \orb(A), O_B \in \orb(B)$. 
%	Let us write $n_A = |\supp(x_{O_A})|$ and $n_B = |\supp(x_{O_B})|$, with further calculations assuming $n_B \le n_A$ (the calculations are similar for $n_A \le n_B$).
%	Then, we have:
%	% for any $X_A \in O_A,X_B \in O_B$ and $\pi' \in \fix(\Atoms \setminus \supp(X_A, X_B))$ we have:
%	\begin{align*}
%		|F(O_A, O_B)| &= |\bigcup_{\substack{\vec b \lesssim \vec b_0,\\ \{\vec a\} \subseteq \supp(x_{O_A})}}\orbs((x_{O_A}, \move{\vec b}{\vec a}\cdot x_{O_B}))|\\
%		&=
%		|\{\move{\vec b}{\vec a} \mid \{\vec a\} \subseteq \supp(X_A), \vec b \lesssim \vec b_0\}|
%		\\&= |\{\sigma: \supp(X_B) \rightharpoonup \supp(X_A) \mid \sigma\text{ is a 1-1 function}\}|
%		\\ &= n_A(n_A - 1)\dots(n_A - n_B + 1) 
%		\\&+ (n_B)(n_A)(n_A - 1)\dots(n_A - n_B + 2) + \dots + 1
%		\\ &=\sum^{n_B}_{i = 0} \choice{n_B}{i} n_A(n_A - 1)\dots(n_A - i + 1)  
%		\\&\le         \sum^{n_B}_{i = 0} 
%		\choice{n_B}{i}
%		(n_A-1)^i \cdot \frac{n_A}{n_A - 1}
%		\\ &\le {n_A}^{n_B} \cdot \frac{n_A}{n_A-1} \le {n_A}^{n_B}(1 + o(1)) 
%	\end{align*}
%	Let us set $n_Y = \regindex(Y)$ for $Y \in \{A, B\}$, and let $n_1 = \max(\regindex(A),\regindex(B)), n_2 = \min(\regindex(A),\regindex(B))$.
%	Then we have:
%	\[
%	|\orb(A \times B)| \le |\orb(A)|\cdot|\orb(B)|\cdot {n_1}^{n_2}(1 + o(1))
%	\]
%\end{proof}

%\begin{lemma}\label{lem:orbtoorbclos}
%	Let $\Xcal$ be a nominal set and $X,Y \subseteq \Xcal$. 
%	Then:
%	\begin{enumerate}
%	\item\label{lem:orbtoorbclos:1} $|\orb(\cls{}(X))| \le |\orb[\supp(X)](X)|$,
%     \item\label{lem:orbtoorbclos:2} if $X$ has empty support then $\cls{}(X \times Y) = X \times \cls{}(Y)$,
%     \end{enumerate}
%\end{lemma}
%
%\begin{proof}
%	For 1, we define the function $f: \orb(\cls{}(X)) \to \orb[\supp(X)](X)$ as follows:
%	\[
%	f(O) = \orbs[\supp(X)](y) \text{ for some } y \in O \cap X
%	\]
%	Let us choose $O_1, O_2 \in \orb(\cls{}(X))$ such that $O_1 \neq O_2$ and $\hat x_1 \in O_1 \cap X, \hat x_2 \in O_2 \cap X$.
%	Suppose $f(O_1) = \orbs[\supp(X)](\hat x_1) = \orbs[\supp(X)](\hat x_2) = f(O_2)$. 
%	This means that $\hat x_1, \hat x_2 \in \orbs[\supp(X)](\hat x_1)$, 
%	and so there is some $\pi \in \fix(\supp(X))$ such that $\pi \cdot \hat x_1 = \hat x_2$. 
%	As $\fix(\supp(X)) \subseteq \Perm$, this implies $\hat x_2 \in O_1$, which is a contradiction. 
%	Thus, $f$ is a 1-1 function and $|\orb(\cls{}(X))| \le |\orb[\supp(X)](X)|$ as required.
%	\\
%    For 2, we have that 
%	\begin{align*}
%		\cls{}(X \times Y) &= \{\pi \cdot (x, y) \mid \pi \in \Perm, (x,y) \in X \times Y\}
%		\\ &= \{(\pi \cdot x, \pi \cdot y) \mid \pi \in \Perm, (x, y) \in X \times Y\}
%	\end{align*} 
%	As $\supp(X) = \emptyset$, this means that $X$ is closed under permutation, and so we obtain:
%	\begin{align*}
%		\{(\pi \cdot x, \pi \cdot y) \mid \pi \in \Perm, (x, y) \in X \times Y\} &= \{(x, \pi \cdot y) \mid \pi \in \Perm, (x, y) \in X \times Y\}\\
%		&= X \times \{\pi \cdot y \mid \pi \in \Perm, y \in Y\} \\ &= X\times \cls{}(Y)
%	\end{align*}
%	\end{proof}




In the following proofs, when discussing any sequence of names $\vec a$, it is assumed that all names in $\vec a$ are distinct.


% \subsection{Proof of \cref{lem:sizebound}}
\begin{lemma}\label{lem:sizebound}
	For every formula $\phi$ such that for every subformula $(\mu X(\vec x).\phi')(\vec u)$ of $\phi$ we have $X$ is free in $\phi'$, then $|\supp(\phi)|+2\|\phi\|+\zeta(\phi)\leq|\phi|$ where $\zeta(\phi)$ is the number of free name variables (of the form $x, y, \dots$) in $\phi$.
\end{lemma}
\begin{proof}
	We perform an induction on $\phi$ with the following base cases:
	\begin{itemize}
		\item $u=v$: we have that $|\supp(\phi)| + \zeta(\phi) = 2$, $2\|\phi\| = 0$  and $|\phi| = 2$. 
		\item $X(\vec u)$: we have that $|\supp(X(\vec u))| + \zeta(\phi) = |\vec u|$, $2\|X(\vec u)\| = 0$, $\zeta(\phi)=0$ and $|X(\vec u)| = 1 + |\vec u|$.
	\end{itemize}
	Next we perform inductive steps:
	\begin{itemize}
		\item $\phi_1 \lor \phi_2$: 
		% for all $i \in \{1,2\}$, 
		we have that 
		% $\supp(\phi_i) \subseteq |\supp(\phi)|$ and 
		$\|\phi\| = \max(\|\phi_1\|, \|\phi_2\|)$,
		and by the inductive hypothesis we have $|\supp(\phi_i)| + 2\|\phi_i\| + \zeta(\phi_i) \le |\phi_i|$.
		Additionally, we know $\supp(\phi_i) \subseteq \supp(\phi)$ and $\zeta(\phi) \le \zeta(\phi_1) + \zeta(\phi_2)$. 
		By definition, we know $|\phi| = 1 + |\phi_1| + |\phi_2|$, 
		thus, we have:
		\begin{align*}
			|\supp(\phi)| + 2\|\phi\| + \zeta(\phi)&\le (|\supp(\phi_1)| + |\supp(\phi_2)|) + 2(\|\phi_1\| + \|\phi_2\|) + \zeta(\phi_1) + \zeta(\phi_2) 
			\\&\le 1 + |\phi_1| + |\phi_2| = |\phi|
		\end{align*}
		%
		\item $\lnot \phi'$: we have that $\|\lnot \phi'\| = \|\phi'\|$, and by the inductive hypothesis we have $|\supp(\phi')| + 2\|\phi'\| + \zeta(\phi') \le |\phi'|$. 
		Additionally, we have that $\supp(\phi') = \supp(\lnot \phi')$ and $\zeta(\phi) = \zeta(\phi')$.
		By definition, we have that $|\lnot \phi'| = 1 + |\phi'|$, thus we obtain:
		\[
		|\supp(\lnot \phi')| + 2\|\lnot \phi'\| + \zeta(\phi)= |\supp(\phi')| + 2\|\phi'\| + \zeta(\phi')\le 1 + |\phi'|
		\]
		%
		\item $\diam{\ell} \phi'$: 
		we have that $\|\diam{\ell}\phi'\| = \|\phi'\|$, and by the inductive hypothesis we have $|\supp(\phi')| + 2\|\phi'\| + \zeta(\phi') \le |\phi'|$. 
		Additionally, as $\phi$ could have more names than $\phi'$ by up to $|\vec u|$, we have that $|\supp(\phi)| \le |\supp(\phi')| + |\vec u|$, and as no variables become free we have $\zeta(\phi) = \zeta(\phi')$.
		By definition, we have that $|\diam{t, \vec u} \phi'| = 1 + |\vec u| + |\phi'|$, thus we obtain:
		\begin{align*}
			|\supp(\diam{t, \vec u} \phi')| + 2\|\diam{t, \vec u} \phi'\| + \zeta(\phi) 
			&= |\supp(\phi')| + 2\|\phi'\| + \zeta(\phi') + |\vec u|\\
			&\le 1 + |\phi'| + |\vec u| = |\phi|
		\end{align*}
		Similar can be shown for the case of $\fresh{x}\phi'$.
		%
		\item $\bigvee\nolimits_x \phi'$: we have that $\|\bigvee\nolimits_x \phi'\| = 1 + \|\phi'\|$, and by the inductive hypothesis we have $|\supp(\phi')| + 2\|\phi'\| + \zeta(\phi') \le |\phi'|$.
		Additionally, we have that $\supp(\bigvee\nolimits_x \phi') = \supp(\phi')$ and as $x$ is free in $\phi'$ (but not in $\phi$) we have $\zeta(\phi) = \zeta(\phi') - 1$.
		By definition, we have that $|\bigvee\nolimits_x \phi'| = 2 + |\phi'|$, thus we obtain:
		\begin{align*}
			|\supp(\bigvee\nolimits_x \phi')| + 2\|\bigvee\nolimits_x  \phi'\| + \zeta(\phi) &= |\supp(\phi')| + 2(\|\phi'\| + 1) + (\zeta(\phi') - 1)\\
			&= |\supp(\phi')| + 2\|\phi'\| + \zeta(\phi') + 1\\
			&\le 1 + |\phi'| = |\phi|  
		\end{align*}
		%
		\item $(\mu X(\vec x).\phi')(\vec u)$: we have that $\|(\mu X(\vec x).\phi')(\vec u)\| = \|\phi'\| + n$ where $n = |\vec u| = |\vec x|$, and by the inductive hypothesis we have that $|\supp(\phi')| + 2\|\phi'\| + \zeta(\phi') \le |\phi'|$. 
		Additionally, $\phi$ will have at most $n$ many more names than $\phi'$, and so $|\supp(\phi)| \le |\supp(\phi')| + n$. 
		% we have that $\supp((\mu X(\vec x).\phi')(\vec u)) = \supp(\phi')$
		and as each $x \in \{\vec x\}$ is free in in $\phi'$ (but not in $\phi$), then $\zeta(\phi) = \zeta(\phi') - n$.
		By definition we have that $|(\mu X(\vec x).\phi')(\vec u)| = 1 + |\phi'| + 2n$, thus we obtain:
		\begin{align*}
			|\supp((\mu X(\vec x).\phi')(\vec u))| + 2\|(\mu X(\vec x).\phi')(\vec u)\| + \zeta(\phi)
			&\le (|\supp(\phi')| + n) + 2(\|\phi'\| + n) + \zeta(\phi') - n\\
			% &= |\supp(\phi')| + 2\|\phi'\| + \zeta(\phi') + |\vec a|\\
			&\le |\phi'| + 2n + 1 = |\phi|\ 
		\end{align*}
	\end{itemize}
\end{proof}


\begin{lemma}\label{lem:genPartialPermF}
	Let $\phi_0$ be a \fhml\ formula, $\vec a_0, \vec b_0 \in \Atoms^{n}$ where $|\vec a_0| = |\vec b_0| = n$ and let $(\phi, \gamma_1, \theta_1), (\psi, \gamma_2, \theta_2) \in \clos(\phi_0)$.
%	Let 
	Then, the time complexity of finding $\move{\vec a}{\vec b}$ 
	such that $\move{\vec a}{\vec b} \cdot \phi\{\theta_1\}\{\gamma_1\} = \psi\{\theta_2\}\{\gamma_2\}$, $\vec a_0$ is a subsequence of $\vec a$ and $\vec b_0$ is a subsequence of $\vec b$ is bounded by $O(|\phi_0)|)$.
%	\ntnote{I think that everywhere here you assume that sequences like $\vec a,\vec b$ have distinct names?}
\end{lemma}

\begin{proof}
	Let us define the function $\Gamma$ that takes some $(\phi, \theta_1, \gamma_1, \psi, \theta_2, \gamma_2, \vec a, \vec b)$ as parameters and returns some $(\vec a', \vec b')$ such that:
	\[
	\begin{aligned}
		\vec a' &= \vec a, \vec a'' &[\text{s.t.} \{\vec a\} \cap \{\vec a''\} = \emptyset]\\
		\vec b' &= \vec b, \vec b'' &[\text{s.t.} \{\vec b\} \cap \{\vec b''\} = \emptyset]\\
	\end{aligned}
	\]
	\vspace{-1.5ex}
	\[
	\move{\vec a'}{\vec b'} \cdot \phi'\{\theta_1\}\{\gamma_1\} = \psi\{\theta_2\}\{\gamma_2\}
	\]
	or returns $\NO$ if not possible. 
	Moreover, we define the function $\update$ that takes some $(\vec a, \vec b, \vec c, \vec d)$, such that $|\vec a| = |\vec b|$ and:
	\begin{itemize}
		\item if $|\vec c| \neq |\vec d|$ then return $\NO$,
		%
		\item if $\vec c = \varepsilon$ and $\vec d = \varepsilon$, then return $(\vec a, \vec b)$,
		%
		\item if $c_1 \in \{\vec a\}$ and $d_1 \in \{\vec b\}$, then 
		we
		return the result of $$\update(\vec a, \vec b, \vec{c}_{2:}, \vec{d}_{2:})$$ (where for any $\vec x$, $\vec{x}_{2:}$ is $\vec{x}$ excluding the first term) if $\move{\vec a}{\vec b}(c_1) = d_1$, and returns $\NO$ otherwise,
		% 
		\item if $c_1 \notin \{\vec a\}$ and $d_1 \notin \{\vec b\}$, then we 
		return the result of $$\update(\vec a(c_1), \vec b(d_1), \vec{c}_{2:}, \vec{d}_{2:})$$
		% 
		\item if $c_1 \in \{\vec a\}$ and $d_1 \notin \{\vec b\}$ (or vice versa), then we return $\NO$.
	\end{itemize}
	We construct $\Gamma(\phi, \theta_1, \gamma_1, \psi, \theta_2, \gamma_2, \vec a, \vec b)$ inductively on $\phi, \psi$ with the following base cases:
	\begin{itemize}
		\item $\phi = (u_1=u_2)$: 
		if $\psi$ is not of the form $(v_1 = v_2)$, then we return $\NO$. 
		Else, we return $\update(\vec a, \vec b, (u_1,u_2)\{\gamma_1\}, (v_1,v_2)\{\gamma_2\})$.
		% 
		\item $\phi = X(\vec u)$: there must be a unique $\theta \in \theta_1$  such that $X(\vec c)\in\dom(\theta)$, where $\vec c=\vec u\{\gamma_1\}$.
		There are two valid possibilities:
		\begin{enumerate}
			\item $\psi$ is some $X(\vec v)$.
			% 
			\item $\psi$ is some $(\sigma X(\vec x).\psi')(\vec v)$ and $\theta(X(\vec c)) = (\sigma X(\vec x).\phi')(\vec c)$ for $\sigma \in \{\mu, \nu\}$.
		\end{enumerate}
		In either case, 
		let $\vec y_1$ be an ordering of the free variables of $\phi'$, and $\vec y_2$ be an ordering of the free variables of $\psi'$. 
		We return $\update(\vec a, \vec b, (\vec y_1\{\gamma_1\}, \vec c), (\vec y_2, \vec v)\{\gamma_2\})$.
		This base case is similar for when $\psi$ is of the form $X(\vec v)$.
	\end{itemize}
	Next we perform inductive steps:
	\begin{itemize}
		\item if $\phi = \phi_1 \lor \phi_2$ and $\psi \neq \psi_1 \lor \psi_2$, then we return $\NO$. 
		Otherwise, by the inductive hypothesis, we check if $\Gamma(\phi_1, \theta_1, \gamma_1, \psi_2, \theta_2, \gamma_2, \vec a, \vec b)$ is some $(\vec a', \vec b')$. 
		If not, then we return $\NO$.
		If so, by the inductive hypothesis, we return $\Gamma(\phi_2, \theta_1, \gamma_1,\psi_2, \theta_2, \gamma_1, \vec a', \vec b')$.
		Similar can be shown for when $\phi' = \phi_1 \land \phi_2$.
		% 
		\item if $\phi = \diam{t, \vec c}\phi'$ and $\psi\neq \diam{t, \vec d}\psi'$ then we return $\NO$. 
		Else, we check the result of $\update(\vec a, \vec b, \vec c\{\gamma_1\}, \vec d\{\gamma_2\})$.
		If the result is $\NO$, then we return $\NO$. 
		Otherwise if the result is some $(\vec a', \vec b')$, then by the inductive hypothesis, we return $ \Gamma(\phi', \theta_1, \gamma_1, \psi', \theta_2,\gamma_2, \vec a', \vec b')$.
		The case where $\phi = \sq{t, \vec c}\phi'$ is similar.
		% 
		\item if $\phi = \bigvee_x \phi'$ and $\psi \neq \bigvee_x \psi'$, then we return $\NO$. 
		Else, we select some $c \notin \{\vec a\} \cup \{\vec b\} \cup \supp(\phi, \psi)$ and, by the inductive hypothesis,
		we examine the result of $$\Gamma(\phi', \theta_1, (\gamma_1, \{c/x\}), \psi', \theta_2, (\gamma_2, \{c/x\}), (\vec a, c), (\vec b, c))$$
		If this is $\NO$, then we return $\NO$. 
		Else if the result is some $(\vec a', \vec b')$, then we return $(\vec a'', \vec b'')$, where $\vec a''$ is $\vec a'$ excluding $c$, and similarly for $\vec b''$.
		The case where $\phi = \bigwedge_x \phi'$ and $\phi = \fresh{x}\phi'$ is similar.
		% 
		\item if $\phi = (\mu X(\vec x).\phi')(\vec u)$ and $\psi \neq (\mu X(\vec x).\psi')(\vec v)$, then we return $\NO$.
		Otherwise, we check the result of $\update(\vec a, \vec b, \vec u\{\gamma_1\}, \vec v\{\gamma_2\})$. If this is $\NO$, then we return $\NO$. 
		Otherwise if this is some $(\vec a', \vec b')$, then 
		by the inductive hypothesis, we return $$\Gamma(\phi', \theta_1, (\gamma_1, \{\vec u\{\gamma_1\} / \vec x\}), \psi, \theta_2, (\gamma_2, \{\vec v\{\gamma_2\} / \vec x\}), \vec a', \vec b')$$
		Note the case where $\psi$ is of the form $X(\vec v)$ is covered by the base case.
		The case where $\phi = (\nu X(\vec x).\phi')(\vec u)$ is similar.
	\end{itemize}
	For any $(\phi, \gamma_1, \theta_1), (\psi, \gamma_2, \theta_2) \in \clos(\phi_0)$, to determine if there exists a $\vec a, \vec b$ such that $\move{\vec a}{\vec b}\cdot \phi\{\theta_1\}\{\gamma_1\} = \psi\{\theta_2\}\{\gamma_2\}$, it suffices to check if $\Gamma = \Gamma(\phi, \theta_1, \gamma_1, \psi, \theta_2, \gamma_2, \vec a_0, \vec b_0)$ is some $(\vec a', \vec b')$.
	It remains to determine the time complexity of constructing $\Gamma$.  
	The number of inductive calls to $\Gamma$ is bounded by $\min(|\phi|, |\psi|)$, and each such inductive call will internally call $\update$ at most once.
	The complexity of any call $\update(\vec a', \vec b', \vec c, \vec d) = (\vec a'', \vec b'')$ will be $O(|\vec a''| - |\vec a'|)$.  
	Over the course of constructing $\Gamma$, this means the cumulative complexity of all calls to $\update$ will be in $O(\min(|\supp(\phi)| + \|\phi\|, |\supp(\psi)| + \|\psi\|))$
	Thus, the time complexity of constructing $\Gamma$ is:
	\begin{align*}
		&O(\min(|\phi|, |\psi|) + \min(|\supp(\phi)| + \|\phi\|, |\supp(\psi)| + \|\psi\|)) \\
		\le\ &O(|\phi_0| + \supp(\phi_0) + \|\phi_0\|) \\
		\overset{\text{Lemma~\ref{lem:sizebound}}}{\le}\ &O(|\phi_0| + |\phi_0|) \le O(|\phi_0|)
	\end{align*} 
\end{proof}

\subsection*{Proof of \cref{lem:orbitGameSize}(2)}\label{app:complexityboundedgame}
\begin{lemma}
  Given a setup $(\Lcal,\phi_0,s_0,H_0)$ of grade $N$ and $\Gcal=\Gcal_N(\Lcal, \phi_0, s_0, H_0)$,
 the time complexity of constructing $\gorb(\Gcal)$ is in
 	\begin{align*}
	&O(
 	(|\phi_0| + \Phi_\Lcal(|\supp(\phi_0)| + \|\phi_0\|) + \S)
 	\cdot
 	\max(N + 2, |\Scal|
 	) \\&\ \cdot 
 	(|\orb(\Scal)| \cdot|\phi_0|\cdot \frac{(|\supp(\phi_0)|+\|\phi_0\|)!}{|\supp(\phi_0)|!}\cdot (\S + 1)^{|\supp(\phi_0)| + \|\phi_0\| + 1})
 	^2)
 \end{align*}
  \end{lemma}
\begin{proof}
%	We begin by defining the function $T$, which takes as parameter $((\phi', \gamma', \theta'), s', H')$ and returns $\orbs(\phi'\{\theta'\}\{\gamma'\}, s', H')$, so $(\phi', \gamma', \theta') \in \clos(\phi_0)$.
	We begin by defining the set $T$ of orbit representatives of $\pos(\Gcal)$, that is, a triple $((\phi, \gamma, \theta), s, H)$ represents $\orbs(\phi\{\theta\}\{\gamma\}, s, H)$.
	We build the game inductively on $\pos(\Gcal)$, with the base case being that $\orbs(\phi_0, s_0, H_0') \in \pos(\Gcal)$ for some $H_0'\wellbound{N}(s_0,\phi_0,H_0)$, and initially $T = \{((\phi_0, \varepsilon, \varepsilon), s_0, H_0')\}$.
	
	Suppose we are examining a position $\orbs( \phi,  s,  H)$, 
	with representative $((\phi_i, \gamma_i, \theta_i), s_i, H_i)$, so $(\phi_i, \gamma_i, \theta_i), s_i, H_i) \in T$. 
	For all moves $(\phi_i\{\theta_i\}\{\gamma_i\}, s_i, H_i) \to (\phi'\{\theta'\}\{\gamma'\}, s', H')$ such that $(\phi', \gamma', \theta') \in \clos(\phi_0)$, we need to check if $\orbs(\phi'\{\theta'\}\{\gamma'\}, s', H') \in \pos(\Gcal)$. 
%	By \lemmapoint{lem:freshClosProps}{2}, there exists a triple $(\phi', \gamma', \theta') \in \clos(\phi_0)$ such that $\phi'\{\theta'\}\{\gamma'\} = \phi$, and we assume this triple can be found in constant time.
%	\ntnote{\scriptsize there is no guarantee that is can be found efficiently. But, instead of checking the transitions of $(\phi_i\{\theta_i\}\{\gamma_i\}, s_i, H_i) \to (\phi, s', H')$ we can look directly at transitions $(\phi_i\{\theta_i\}\{\gamma_i\}, s_i, H_i) \to (\phi'\{\theta'\}\{\gamma'\}, s', H')$ where $(\phi',\gamma',\theta')$ is a child of $(\phi_i,\gamma_i,\theta_i)$ in the clos construction. This can be done in constant time}
	Thus, we need to check if for any $((\phi_1, \gamma_1, \theta_1), s_1, H_1) \in T$:
	$$\pi \cdot ((\phi', \gamma', \theta'), s', H') = ((\phi_1, \gamma_1, \theta_1), s_1, H_1)$$ 
	for some permutation $\pi$ (i.e., they are nominally equivalent).
%	This means we must see if there is a permutation $\pi \in \Perm$ such that $\pi \cdot ((\phi', \gamma', \theta'), s', H') = ((\phi_1,\gamma_1, \theta_1), s_1, H_1)$.
	For such a $\pi$ to exist, we must satisfy the following conditions:
	\begin{enumerate}
		\item $|H'| = |H_1|$.
		%
		\item There must exist a minimal $\move{\vec a}{\vec b}$ such that $\move{\vec a}{\vec b} \cdot (\phi', \gamma', \theta') = (\phi_1, \gamma_1, \theta_1)$; 
	 	by Lemma~\ref{lem:genPartialPermF}, this can be constructed in $O(|\phi_0|)$.
		%
		\item From $\vec a, \vec b$ as in the previous condition, the oracle $\Psi_\Lcal(\vec a, \vec b, s_1, s')$ must find a permutation.
		Each such check is in $O(\Phi_\Lcal(|\vec a|)) \le O(\Phi_\Lcal(|\supp(\phi)| + \|\phi\|)) \le O(\Phi_\Lcal(|\supp(\phi_0)| + \|\phi_0\|))$. 
%		which by Lemma~\ref{lem:sizebound}, is in $O(\Phi_\Lcal(|\phi_0|))$
		%
		\item The number of redundant names in the histories must be the same, that is, $|H' \setminus \supp(\phi, s')| = |H_1 \setminus \supp(\phi_1\{\theta\}\{\gamma_1\}, s_1)|$. 
		Each such check is in $O(|\supp(\phi_0)| + \|\phi_0\| + \regindex(\Scal)) = O(N)$.
	\end{enumerate}
%	e require that $|H'| = |H_1|$, otherwise there is no such permutation $\pi$.
%	 We can first construct a minimal permutation $\move{\vec a}{\vec b}$ such that $\move{\vec a}{\vec b} \cdot (\phi', \gamma', \theta') = (\phi_1, \gamma_1, \theta_1)$; 
%	 by Lemma~\ref{lem:genPartialPermF}, this can be constructed in $O(|\phi_0|)$.
%	If there is no such permutation, then $((\phi', \gamma', \theta'), s', H') $ is not nominally equivalent to $((\phi_1,\gamma_1, \theta_1), s_1, H_1)$.
%	Otherwise, the oracle $\Psi_\Lcal(\vec a, \vec b, (q_1, \rho_1), s')$ can attempt to find a permutation $\pi'$. 
%	If $\Psi_\Lcal(\vec a, \vec b, (q_1, \rho_1), s')$ returns $\NO$, then there is no such $\pi$;
%	otherwise, it remains to examine $H', H_1$.
%	We have that $\move{\vec a}{\vec b}\cdot\supp(\phi, s') = \supp(\phi_1\{\theta^*\}\{\gamma_1\}, s_1)$, and so it suffices to check if both $H'$ and $H_1$ have the same number of redundant names.
%	Thus, we check $|H' \setminus \supp(\phi, \rho')| = |H_1 \setminus \supp(\phi_1\{\theta^*\}\{\gamma_1\}, \rho_1)|$ where each such check is in $O(|\supp(\phi_0)| + \|\phi_0\| + \regindex(\Scal))$.
%	If this is not the case, then there is no such permutation.
	Each such nominal equivalence check is in:
	\[
	O(
	|\phi_0| + \Phi_\Lcal(|\supp(\phi_0)| + \|\phi_0\|) + N
	)
	\]
%		&O(|\phi_0| + |\supp(\phi_0)| + \|\phi_0\| + \regindex(\Scal) + \Phi_\Lcal(|\vec a|, |s'|, |(q_1, \rho_1)|)) \\\overset{\text{Lemma~\ref{lem:sizebound}}}{=}\ & O(2|\phi_0| + \regindex(\Scal) + \Phi_\Lcal(\min(|\supp(\phi')|, |\supp(\phi_1)|), \max(\Lcal), \max(\Lcal))
%		\end{align*}
	If there is no such $((\phi_1, \gamma_1, \theta_1), s_1, H_1) \in T$ that satisfies all above conditions for $((\phi', \gamma', \theta'), s', H')$, then we add $(\phi', \gamma', \theta'), s', H')$ to $T$.
%	\[
%		T((\phi', \gamma', \theta'), s', H')) = \orbs(\phi, s', H')
%	\]
		
	The moves that can be made from each $\orbs(\phi, s, H)$ would need to be checked at most once, thus the maximum number of positions that will need to be checked is bounded by the maximum number of moves that reach unique targets.
%	\ntnote{\tiny not sure about unique targets. Note that, somewhere in this calculation, the number of labels used in the game should appear, which are the number of tags times the names they can use ($N+2$) to the power of the tag's arity}
	This can be quantified as $\max(\supp(\phi, s) + 2, \ndbb(\Lcal))$, that is:
	\begin{enumerate}
		\item The maximum number of unique names in the support of $\phi, s$, plus one additional name that could be in the history, and one fresh name.
		\item The maximum number of unique moves at any given position.
	\end{enumerate}
	The number of elements that a nominal equivalence check required for each such target is bounded by $|\pos(\Gcal)|$.
	Finally, the above process is done for every position in $\pos(\Gcal)|$, 
	 thus the final complexity is in:
	\begin{align*}
 		&O(
 		(|\phi_0| + \Phi_\Lcal(|\supp(\phi_0)| + \|\phi_0\|) + \S) 			\cdot
 			\max(\supp(\phi, s) + 2, \ndbb(\Lcal)
 		)\cdot |\pos(\Gcal)|^2)
 		\\
 		\le\ &O(
 		(|\phi_0| + \Phi_\Lcal(|\supp(\phi_0)| + \|\phi_0\|) + \S)
 		\cdot
 		\max(N + 2, |\Scal|
 		) \\&\ \cdot 
 		(|\orb(\Scal)| \cdot|\phi_0|\cdot \frac{(|\supp(\phi_0)|+\|\phi_0\|)!}{|\supp(\phi_0)|!}\cdot (\S + 1)^{|\supp(\phi_0)| + \|\phi_0\| + 1}\cdot(1+o(1)))
 		^2)
% 		\\
% 		\overset{\text{Lemma~\ref{lem:sizebound}}}{\le}\ &O(
% 		(|\phi_0| + \Phi_\Lcal(|\phi_0|) + \S)
% 		\cdot
% 		\max(|\phi_0| + \regindex(\Scal) + 2, |\fstates|
% 		)
% 		\\
% 		=\ &O(
% 		(|\phi_0| + \Phi_\Lcal(|\phi_0|) + \S)
% 		\cdot
% 		\max(\S + 2, |\fstates|
% 		)
 	\end{align*}
 \end{proof}

\subsection{Proof of \cref{thm:MCinFHML}}
\MCinFHML*
  \begin{proof}
To verify whether $(s, H) \in \fstates$ satisfies $\phi_0$, we must check whether $(s, H) \in \sem[\xi_0]{\phi_0}$ where $\xi_0$ is an empty $\Ucal$-variable assignment. 
First, we set $\phi = !(\phi_0)$, which (by Lemma~\ref{lem:negfreesize}) will have size bounded by $|\phi_0|$. 
From this, we can construct the orbits parity game $\gorb(\Gcal_\S)$ of $\Gcal_\S = \Gcal_\S(\Lcal, \phi, s, H)$ where $\S = |\phi_0| + \|\phi_0\| + \regindex(\Scal)$, 
By Theorem~\ref{thm:FnomSatPar}, we have that if $(s, H) \in \sem[\xi_0]{\phi_0}$ then defender has a winning strategy from $\fgame = \fgame(\Lcal, \phi, s, H)$.
Using Lemma~\ref{lem:boundingisnombisim}, the relation $\Rcal_{(\Lcal, \phi, s, H)} \subseteq \pos(\fgame) \times \pos(\Gcal_\S)$ is a nom-bisimulation, 
and by Lemma~\ref{lem:nombisimequiv} then defender wins in $\Gcal$ from $(s, H, \xi_0, \phi)$ iff they win in $\Gcal_\S$ from $(s, H, \xi_0, \phi)$.
Similarly, using Lemmaa~\ref{lem:parityBisim} and Corollary~\ref{cor:winningNom} then defender wins from $\Gcal_\S$ from $(s, H, \xi_0, \phi)$ iff defender wins in $\gorb(\Gcal_\S)$ from $\orbs(s, H, \xi_0, \phi)$. 
%and by Lemma~\ref{lem:bisimfromnomeq} we know that $\Rcal' = \{(\orb(P_1), \orb(P_2))\mid (P_1,P_2) \in \Rcal\}$ is a bisimulation between $\gorb(\fgame)$ and $\gorb(\Gcal_\S)$.
%Using Lemma~\ref{lem:parityBisim} we have that $\fgame$ is bisimilar to $\gorb(\fgame)$, by Lemma~\ref{lem:parityBisimWinning} if defender has a winning strategy from $\fgame$ then defender must also have a winning strategy in $\gorb(\fgame)$, and using the same lemma, the defender would hence also have a winning strategy in $\gorb(\Gcal_\S)$. 
% Using Lemma~\ref{lem:parityBisim} and Lemma~\ref{lem:parityBisimWinning}, if defender has a winning strategy from $\Gcal$ then defender must also have a winning strategy in $\Gcal'$, hence it suffices to find the winning regions of $\Gcal'$. 
% Let us set 
% % $n_{\Ccal_\Acal} = \max\{(|\supp(s)|) \mid s \in \Ccal_\Acal\}$, 
% $n_{\phi_0} = |\supp(\phi_0)| + \|\phi_0\|$, and let $n_1 = \max(\kappa, n_\phi_0), n_2 = \min(\kappa, n_\phi_0)$.
Then, by Lemma~\ref{lem:fnomparsize}, $\gorb(\Gcal_\S)$ can be constructed in:
    \begin{align*}
	&O(
	(|\phi_0| + \Phi_\Lcal(|\supp(\phi_0)| + \|\phi_0\|) + \S)
	\cdot
	\max(N + 2, |\Scal|
	) \\&\ \cdot 
	(|\orb(\Scal)| \cdot|\phi_0|\cdot \frac{(|\supp(\phi_0)|+\|\phi_0\|)!}{|\supp(\phi_0)|!}\cdot (\S + 1)^{|\supp(\phi_0)| + \|\phi_0\| + 1}\cdot(1+o(1)))
	^2)
\end{align*}
%\begin{align*}
%	&O((2|\phi| + \Phi_\Lcal(|\phi|) + \regindex(\Scal))\cdot\max(|\phi| + \regindex(\Scal) + 2, |\fstates|))
%	\\
%	\le\ &O((2|\phi_0| + \Phi_\Lcal(|\phi_0|) + \regindex(\Scal))\cdot\max(|\phi_0| + \regindex(\Scal) + 2, |\fstates|))
%\end{align*}
and has size:
\[
n \cdot|\phi_0|\cdot \frac{(|\supp(\phi_0)|+\|\phi_0\|)!}{|\supp(\phi_0)|!}\cdot (\S + 1)^{|\supp(\phi_0)| + \|\phi_0\| + 1}\cdot(1+o(1))
\]
The winning regions of $\gorb(\Gcal_\S)$ can be calculated in time bounded by (cf.
~\cite{CaludeJKLS22}):
%\begin{align*}
\[
O((n \cdot|\phi_0|\cdot \frac{(|\supp(\phi_0)| + \|\phi_0\|)!}{|\supp(\phi_0)|!}\cdot (\S + 1)^{|\supp(\phi_0)| + \|\phi_0\| + 1})^{\log(d) + 6})
\]
% &O((n \cdot|\phi_0|\cdot \frac{(|\phi_0|)!}{|\supp(\phi_0)|!}\cdot (\S + 1)^{|\phi_0| + 1}(1 + o(1)))^{\log(d) + 6}\\
% &\ + (2|\phi_0| + \Phi_\Lcal(|\phi_0|) + \regindex(\Scal))\cdot\max(|\phi_0| + \regindex(\Scal) + 2, |\fstates|))
% \\
% \le\ &O((n \cdot|\phi_0|\cdot \frac{(|\phi_0|)!}{|\supp(\phi_0)|!}\cdot (\S + 1)^{|\phi_0| + 1}(1 + \epsilon))^{\log(d) + 6} 
% \\&\ + (2|\phi_0| + \Phi_\Lcal(|\phi_0|) + \regindex(\Scal))\cdot\max(|\phi_0| + \regindex(\Scal) + 2, |\fstates|))
%\end{align*}

And so, the final complexity is in:
\begin{align*}
	&O((n \cdot|\phi_0|\cdot \frac{(|\supp(\phi_0)| + \|\phi_0\|)!}{|\supp(\phi_0)|!}\cdot (\S + 1)^{|\supp(\phi_0)| + \|\phi_0\| + 1})^{\log(d) + 6}
	\\ &\ + (
	(|\phi_0| + \Phi_\Lcal(|\supp(\phi_0)| + \|\phi_0\|) + \S)
	\cdot
	\max(N + 2, |\Scal|
	) \\&\ \cdot 
	(n \cdot|\phi_0|\cdot \frac{(|\supp(\phi_0)|+\|\phi_0\|)!}{|\supp(\phi_0)|!}\cdot (\S + 1)^{|\supp(\phi_0)| + \|\phi_0\| + 1})
	^2)
	)
	\\
	=\ &O((n \cdot|\phi_0|\cdot \frac{(|\phi_0|)!}{|\supp(\phi_0)|!}\cdot (\S + 1)^{|\phi_0| + 1})^{\log(d) + 6}
	\\ &\ + (
	(|\phi_0| + \Phi_\Lcal(|\phi_0|) + \S)
	\cdot
	\max(N + 2, |\Scal|
	) \\&\ \cdot 
	(n \cdot|\phi_0|\cdot \frac{(|\phi_0|)!}{|\supp(\phi_0)|!}\cdot (\S + 1)^{|\phi_0| + 1})
	^2)
	)
\end{align*}
\end{proof}

 
% \subsection{Proof of~\cref{app:FRAGameSize}}
% \FRAGameSize*
% \begin{proof}
% 	The proof is symmetric to that of Lemma~\ref{lem:orbitGameSize}, with the following change: 
% 	when checking for a $\pi$ such that
% 	$$\pi \cdot ((\phi', \gamma', \theta'), (q', \rho'), H') = ((\phi_1, \gamma_1, \theta_1), (q_1, \rho_1), H_1)$$ 
% 	the conditions differ as such:
% 	\begin{enumerate}
% 		\item $|H'| = |H_1|$.
% 		%
% 		\item $q' = q_1$, if so we can obtain a $\move{\vec a}{\vec b}$ such that $\move{\vec a}{\vec b} = (\rho'^{-1}) \circ \rho_1$; this can be constructed in $O(r)$. 
% 		%		
% 		\item Using $\vec a, \vec b$ as in the previous condition, there must exist a $\move{\vec c}{\vec d}$ such that $\move{\vec c}{\vec d}(\vec a) = \vec b$ and $\move{\vec c}{\vec d} \cdot(\phi', \gamma', \theta') = (\phi_1, \gamma_1, \theta_1)$; 
% 		by Lemma~\ref{lem:genPartialPermF}, this can be constructed in $O(|\phi_0|)$.
% 		%
% 		\item The number of redundant names in the histories must be the same, that is, $|H' \setminus \supp(\phi, \rho')| = |H_1 \setminus \supp(\phi_1\{\theta\}\{\gamma_1\}, \rho_1)|$. 
% 		Each such check is in $O(|\supp(\phi_0)| + \|\phi_0\| + r)$
% %		 which by Lemma~\ref{lem:sizebound}, is in $O(|\phi_0| + r)$.
% 	\end{enumerate}
% 	This makes each such nominal equivalence check in:
% 	\begin{align*}
% 		&O(r + |\phi_0| + |\supp(\phi_0)| + \|\phi_0\| + r)
% 		\\
% 		=\ &O(2r + + |\phi_0| + |\supp(\phi_0)| + \|\phi_0\|)
% 	\end{align*}
% 	Thus, the final complexity is in:
% 	\begin{align*}
% 		&O(
% 		(2r + |\phi_0| + |\supp(\phi_0)| + \|\phi_0\|)
% 		\cdot
% 		\max(\supp(\phi, \rho) + 2, \ndbb(\Lcal))
% 		\\
% 		&\ \cdot 
% 		(|Q| \cdot|\phi_0|\cdot \frac{(|\supp(\phi_0)|+\|\phi_0\|)!}{|\supp(\phi_0)|!}\cdot (\S + 1)^{|\supp(\phi_0)| + \|\phi_0\| + 1}\cdot(1+o(1)))
% 		^2)
% 		\\
% 		\le\ &O(
% 		(2r + |\phi_0| + |\supp(\phi_0)| + \|\phi_0\|)
% 		\cdot
% 		\max(|\supp(\phi_0)| + \|\phi_0\| + r + 2, |Q|
% 		)
% 		\\
% 		&\ \cdot 
% 		(|Q| \cdot|\phi_0|\cdot \frac{(|\supp(\phi_0)|+\|\phi_0\|)!}{|\supp(\phi_0)|!}\cdot (\S + 1)^{|\supp(\phi_0)| + \|\phi_0\| + 1}\cdot(1+o(1)))
% 		^2)
% 	\end{align*}
% \end{proof}
%	 
%	
%	Each $\orbs(\phi, s, H)$ would need to be examined at most once, and the number of moves that would need to be checked is the maximum of the number of moves that reach a unique target, that is, $\max(\supp(\phi, s) + 2, \ndbb(\Lcal))$. \hbnote{+2 should come for two fresh names - one forgotten and one fresh} 
%	If there is no such $((\phi_1, \theta_1, \gamma_1), (q_1, \rho_1), H_1) \in \dom(T)$ such that $((\phi', \theta', \gamma'), s', H') \nomeq ((\phi_1, \theta_1, \gamma_1), s_1, H_1)$, then we add $\orbs(\phi, s', H')$ to $\pos(\Gcal)$ and update $T$ to $T \cup \{((\phi', \gamma', \theta'), s', H'), \orbs(\phi, s', H')\}$. 
%	Thus, the final complexity is in
%	\begin{align*}
%			&O(2|\phi_0| + \regindex(\Scal) + \Phi_\Lcal(\min(|\supp(\phi')|, |\supp(\phi_1)|), \max(\Lcal), \max(\Lcal))  \\&\ \cdot \max(\supp(\phi, s) + 1, \ndbb(\Lcal))
%			\\ 
%			\overset{\text{\lemmapoint{lem:freshClosProps}{5}}}{\le} & O(\max(|\phi_0| + \|\phi_0\| + \max(\Lcal) + 1, |\Scal|)\\
%			&\ \cdot (2|\phi_0| + \regindex(\Scal) + \Phi_\Lcal(|\phi_0| + \|\phi_0\|, \max(\Lcal), \max(\Lcal))))
%		\end{align*}
%\\
%Next, we show the size bound of $\gorb(\Gcal_\S)$.
%	By definition and \lemmapoint{lem:freshClosProps}{2}, we know that $\pos(\Gcal_\S) \subseteq \fstates \times \clos(\phi_0)$.
%	% \fixednote{I thought $\hat{ \cal S}$ already contained the history. $H\in\mathcal{P}(\Atoms^{N+1})$ is a set of $N+1$-tuples of names. Also, the set of positions is not the whole of $\hat{\cal S}\times\clos(\phi_0)$.}
%	Then, by \lemmapoint{lem:orbtoorbclos}{1} and \lemmapoint{lem:orbtoorbclos}{2}, we have that $\pos(\gorb(\Gcal_\S)) \subseteq \orb(\cls{}(\fstates \times \clos^*(\phi_0))) = \orb(\fstates \times \cls{}(\clos^*(\phi_0)))$.
%	For every $(s, H, \phi_0') \in \Gcal_\S$, we have that $\supp(s) \subseteq H$, and so we have $|\orb(\fstates)| \le |\orb(\Scal)| \cdot (\S + 1)$.
%	% \fixednote{the last equality seems wrong. In any case, there is confusion between states and states-with-histories}
%	By Lemma~\ref{lem:cartesianOrbitSize}, we have that 
%	% \fixednote{is there a lemma saying that the O of the Cl is the cal-O-supp (second line)?}
%	\begin{align*}
%		|\orb(\cls{}(\fstates \times \clos^*(\phi_0))| \le& |\orb(\Scal)| \cdot (\S + 1) \cdot|\orb(\cls{}(\clos^*(\phi_0)))|\cdot (\S + 1)^{|\supp(\phi_0)| + \|\phi_0\|}(1 + o(1)) 
%		\\ \overset{\text{Lemma~\ref{lem:orbtoorbclos}(\ref{lem:orbtoorbclos:1})}}{\le}& |\orb(\Scal)| \cdot|\orb[\supp(\phi_0)](\clos^*(\phi_0))|\cdot (\S + 1)^{|\supp(\phi_0)| + \|\phi_0\| + 1}(1 + o(1))
%	\end{align*}
%	Using \lemmapoint{lem:freshClosProps}{3}, we obtain:
%	\begin{align*}
%		&|\orb(\Scal)|\cdot|\orb[\supp(\phi_0)](\clos^*(\phi_0))|\cdot (\S + 1)^{|\supp(\phi_0)| + \|\phi_0\| + 1}(1 + o(1)) \\
%		\le\ & |\orb(\Scal)| \cdot|\phi_0|\cdot \frac{(|\supp(\phi_0)|+\|\phi_0\|)!}{|\supp(\phi_0)|!}\cdot (\S + 1)^{|\supp(\phi_0)| + \|\phi_0\| + 1}(1 + o(1)) \\
%		\overset{\text{Lemma~\ref{lem:sizebound}}}{\le}\ & |\orb(\Scal)| \cdot|\phi_0|\cdot \frac{(|\phi_0|)!}{|\supp(\phi_0)|!}\cdot (\S + 1)^{|\phi_0| + 1}(1 + o(1))
%	\end{align*}  
%\end{proof}



%\begin{lemma}[Restatement of Lemma~\ref{lem:orbitGameSize}]
%	Given $\phi_0,\Lcal,s_0,H_0$ as above, the time complexity of constructing $\gorb(\Gcal_\S(\Lcal, \phi_0, s_0, H_0))$ is in $O(\dots)$. \ntnote{fill this in}
%\end{lemma}	
%
%\begin{proof}
%	To verify whether $((q, \rho), H)$ satisfies $\phi$, we must check whether $((q, \rho), H) \in \sem[\xi_0]{\phi}$ where $\xi_0$ is an empty $\Ucal$-variable assignment. 
%	First, we set $\hat \phi = !(\phi)$, which (by Lemma~\ref{lem:negSize}) will have size bounded by $|\phi|$. 
%	From this, we can construct the orbits parity game $\gorb(\Gcal_\S)$ of $\Gcal_\S = \Gcal_\S(\Lcal, \phi, (q, \rho), H)$, which by Lemma~\ref{lem:freshOrbitsGameConstructionFRA} can be done in
%	\[
%	O(\max(|\phi| + \|\phi\| + r + 2, |\Scal|) \cdot 2(r + |\phi|))
%	\] 
%	By Theorem~\ref{thm:FnomSatPar}, we have that if $((q, \rho), H) \in \sem[\xi_0]{\phi}$ then defender has a winning strategy from $\fgame = \fgame(\Lcal, \phi, (q,\rho), H)$.
%	Using Lemma~\ref{lem:nomeqbisim}, the relation $\Rcal_{(\Lcal, \phi, (q, \rho), H)} \subseteq \pos(\fgame) \times \pos(\Gcal_\S)$ is a bisimulation up to nominal equivalence, and by Lemma~\ref{lem:bisimfromnomeq} we know that $\Rcal' = \{(\orb(P_1), \orb(P_2))\mid (P_1,P_2) \in \Rcal\}$ is a bisimulation between $\gorb(\fgame)$ and $\gorb(\Gcal_\S)$.
%	Using Lemma~\ref{lem:parityBisim} we have that $\fgame$ is bisimilar to $\gorb(\fgame)$, by Lemma~\ref{lem:parityBisimWinning} if defender has a winning strategy from $\fgame$ then defender must also have a winning strategy in $\gorb(\fgame)$, and using the same lemma, the defender would hence also have a winning strategy in $\gorb(\Gcal_\S)$. 
%	% Using Lemma~\ref{lem:parityBisim} and Lemma~\ref{lem:parityBisimWinning}, if defender has a winning strategy from $\Gcal$ then defender must also have a winning strategy in $\Gcal'$, hence it suffices to find the winning regions of $\Gcal'$. 
%	% Let us set 
%	% % $n_{\Ccal_\Acal} = \max\{(|\supp(s)|) \mid s \in \Ccal_\Acal\}$, 
%	% $n_{\phi} = |\supp(\phi)| + \|\phi\|$, and let $n_1 = \max(\kappa, n_\phi), n_2 = \min(\kappa, n_\phi)$.
%	Then, by Lemma~\ref{lem:fnomparsize}:
%	\[
%	|\gorb(\Gcal_\S)| \le n\cdot(\S + 1)\cdot|\phi|\cdot \frac{(|\supp(\phi)|+\|\phi\|)!}{|\supp(\phi)|!}\cdot (\S + 1)^{|\supp(\phi)| + \|\phi\|}(1 + o(1))
%	\]
%	The winning regions of $\gorb(\Gcal')$ can be calculated in time bounded by (cf.
%	~\cite{CaludeJKLS22}):
%	\[
%	O((n \cdot|\phi|\cdot \frac{(|\supp(\phi)|+\|\phi\|)!}{|\supp(\phi)|!}\cdot (\S + 1)^{|\supp(\phi)| + \|\phi\| + 1}(1 + o(1)))^{\log(d) + 6})
%	\]
%\end{proof}
%%% Local Variables:
%%% mode: latex
%%% TeX-master: "main"
%%% End:
